% Latin 17859, batch 1 — Gallica views 1–20.
% Pass: Opus 5 (claude-opus-5), 2026-09-15 — first pass, unchecked against
% the views by a human. The pass was resumed after an interruption: every
% view had been drafted before the break; this pass checked the drafts
% against the mirror, band by band, re-read the last views (17–20) and
% every doubtful place at full resolution, and assembled the file.
%
% What the views are. Every view is a microfilm image (greyscale, landscape,
% about 7270 × 5840) of the volume lying open: a left page and a right page
% side by side, the gutter near the middle. Each view was split into its two
% halves and read as two pages; in the file each \page{N} holds the left page
% then the right, and a \note{} says where the right page begins.
% View 1 is not transcribed: the inside of the front board with the shelf
% label « LATIN 17,859 », and a guard leaf carrying only the library's
% certificate « Volume de 47 Feuillets. 7 Juin 1870. ». On view 2 the left
% page (back of the guard leaf) is blank. No other page of the batch is blank.
%
% What the twenty views hold. A copy, in one regular italic hand throughout,
% of Maurolico's spherical astronomy, as listed on its own title page.
%
%   v. 2 R    The title page: « Theodosii Sphæricorum Elementorum libri tres.
%             Ex traditione Maurolyci », Menelaus, « Autolyci de sphæra quæ
%             movetur », « Theodosii de Habitationibus », « Euclidis
%             Phoenomena », « Demonstratio et praxis trium tabellarum …
%             sinus recti, Foecundæ, et Beneficæ », « Compendium
%             Mathematicæ », « Maurolyci de sphæra sermo ». Then the source:
%             « Exscripta ex typis excusso codice Messanæ … a Petro Spira
%             mense Augusto Anno M.D.LVIII. », and a provenance note:
%             « Bibliotheca I.U. Thuanorum. mense Iunio. Anno M.DC.XXVIII.
%             favore C.V.D.I.B. Haultini mense decembri mutuo concesso Anno
%             M.DC.XXVII. » An old shelfmark « Nouv. acq. lat. 1099 » in
%             another hand, and a modern pencil note « Ce MS paraît
%             incomplet ».
%   v. 3 L–R  Dedicatory letter « D. Iohanni Cerdæ Methymnensium Duci ac
%             Siciliæ Proregi … Maurolycus S.D. », dated « Messanæ Augusto
%             mense M.D.LVIII. », then « Index lucubrationum Maurolyci »:
%             the editions of others (« Aliena ») and his own works
%             (« Nostra vero sunt »).
%   v. 4 L    Letter « AD Carolum V. Imperatorem maximum Francisci Maurolyci
%             Siculi Abbatis Epistola », dated « Messanæ Iulio mense.
%             M.D.LVI… » (last figure overwritten), and a twelve-line
%             « carmen » to the Emperor in two columns.
%   v. 4 R – 10 L  « Maurolyci Messanensis de Sphæra Sermo »: the sphere
%             as the perfect figure, the roundness and immobility of the
%             earth, the order of the spheres, the circles of the sphere,
%             risings and settings, days and hours, climates, eccentrics and
%             epicycles, latitudes of the planets, trepidation (after
%             Alphonsus), constellations, the division of the houses, and
%             how each astronomical quantity is found by observation.
%             « FINIS. »
%   v. 10 L   Two titles with no text under them: « THEODOSII SPHÆRICORUM
%             ELEMENTORUM EX TRADITIONE MAUROLYCI LIBER PRIMUS. II et III »,
%             « DIFFINITIONES. », with the line « Sed hi collati sunt cum
%             veteri editione Veneta anni 1518. »; then « MENELAI SPHÆRICA EX
%             EIUSDEM TRADITIONE » and « Alibi descripta. » The rest of the
%             page is blank: neither Theodosius nor Menelaus is copied here.
%   v. 10 R – 11 L  Autolycus, « De sphæra quæ movetur », Props. 5–12,
%             without title (see below).
%   v. 11 L – 12 R  « Theodosii De Habitationibus ex traditione Maurolyci.
%             Liber. », Props. 1–12, « FINIS. » Prop. 6 breaks off in its
%             enunciation at the foot of 11 R, and 12 L (the back of the same
%             leaf) opens with Prop. 7.
%   v. 13 L – 15 R  « Euclidis Phænomena ex Traditione Maurolyci »: the
%             preamble, Props. 1–18 with « Additio duarum propositionum »,
%             « Additio trium propositionum » and three « Corollarium »;
%             Maurolico's closing remark on the use of the Sphærica.
%   v. 15 R – 16 R  « Habitationum collatio » (perioeci, antoeci, antipodes,
%             amphiscii, periscii, antiscii, heteroscii) and « DE Astrorum
%             fulsionibus » (the arc of visibility; Venus 5°, the Moon 3°).
%   v. 16 R   « Tabella sinus recti »: sines for 1°–90°, radius 100000,
%             with differences and their sixtieths.
%   v. 17 L   « Tabula foecunda »: tangents (« umbra ») for 1°–90°, with
%             differences and sixtieths, and five values for 89°15′–89°59′.
%   v. 17 R, 18 L, 18 R, 19 L  Problems worked with the tables, numbered
%             1–6: declination and right ascension of a point of the
%             ecliptic, declination and right ascension of a star (Aldebaran,
%             « oculus Tauri », at 2° 50′ Gemini), and the ortive latitude at
%             Messina (38°). See the order of the leaves below.
%   v. 19 R   « Tabella benefica »: secants (« Radius ») for 1°–90°, with
%             differences and sixtieths, and six values for 0° and
%             89°15′–89°59′.
%   v. 20 L   An untitled table of three arcs, « Primus », « Secundus »,
%             « Tertius », for each degree of the ecliptic: the « tabula
%             declinationum et ascensionum » the problems call for (this
%             pass's identification, from its use in Problems 4–5).
%   v. 20 R   The end of a proposition on sines of angles, then Props.
%             XXIII–XXIIII on rectilinear triangles and the enunciation of
%             XXV on the right spherical triangle; view 21 (batch 2) goes on
%             with XXV. This numbering is not the Theodosius above, and is
%             not reconciled here with the Props. XXIII–XXXII of batch 3.
%
% The order of the leaves. Read view by view, the text runs from each right
% page onto the next view's left page (the two sides of one leaf) and from
% each left page onto the right page facing it, except at:
%   – 10 L / 10 R: 10 R opens inside a proposition (« suos seorsum
%     parallelos … describent ») and goes on with Props. 5–9 of a book on
%     the moving sphere. Batch 3, view 48 L, holds « Autolyci de sphæra quæ
%     movetur ex traditione Maurolyci. Liber. » with Props. 1–4, stopping on
%     « et per primam huius puncta singula »: the sentence reads straight on
%     into 10 R. So the Autolycus leaf of batch 3 belongs between views 10 L
%     and 10 R, where the title page puts Autolycus (after Menelaus). This is
%     this pass's reconstruction from the sense; the file keeps the present
%     order.
%   – 17 R, 18 L, 18 R, 19 L: the problems run 1–4 on 18 R (inked 17), on
%     into 19 L (« usque ad dictum latitudinis » / « circulum receptus »),
%     Problem 5 breaks off on « Nam in primis » and continues on 17 R (inked
%     16, « considerandum est »), which runs on into 18 L and Problem 6. The
%     two leaves inked 16 and 17 are bound in the reverse of the text's
%     order. Reconstruction from the sense; the ink numbers run 16, 17 in the
%     present order and the file keeps it.
%   – 17 L / 17 R, 19 R / 20 L / 20 R: tables face text without continuity;
%     20 R opens in a proposition whose beginning is not in the batch.
%
% The numbers on the leaves. Each right page from view 2 to view 20 carries
% at the top right a number in ink, larger than the text: 1, 2, 3, … 19, one
% per leaf, consecutive (view N shows the number N − 1). Left pages carry
% none. The series goes on in batch 3, which reads 40–47 on views 41–48, and
% the guard leaf of view 1 certifies « Volume de 47 Feuillets. 7 Juin
% 1870. »: the ink series and the library's count of 47 leaves agree, which
% suggests that these ink numbers are the foliation the library counted.
% But the figures are ink, not the thin pencil of a library foliation, and
% the stroke could be the copyist's; as in batch 3, no \folio{} is written,
% and each number is recorded in the \note{} opening its right page. The
% question is left to a human reader of the volume. Other numbers: problem
% numbers « 1. »–« 6. » in the margins (views 17–19); « Propositio » with
% its number in the margins (views 10–15), set in bold at the head of the
% enunciation; the proposition numbers cited in the demonstrations.
% A round stamp « Bibliothèque impériale MSS » in the right margin of 2 R
% and of 20 R.
%
% Letters, glyphs, notation. As in batch 3: the long s is set s; the hand's
% one looped form for u and v is set u or v by sound; its initial i, which
% looks like a capital J (« In », « Iam », « Iohannes »), is set I, and
% lower case inside a sentence; « æ » is written as a hooked a/e ligature and
% set æ (the interrupted draft had set it « a » at the end of words; this
% pass restored it, reading each case against the leaf and the grammar). The
% spelling is the copy's: « Ecclipsis », « ecclipticæ », « periferia »
% beside « peripheria », « diffinitio », « autor », « quatuor », « hyems ».
% Contractions are kept as written and expanded in a \note{} at first
% occurrence: « dria » (differentia), « quantite » and « adnotatibus » with a
% stroke, « dispõo », « declinaẽ », « mũdo », « inclinatũ »,
% « horizontib\textsuperscript{9} », « compl\textsuperscript{to} »,
% « septen\textsuperscript{lis} ». In citations the figure 2 is a short
% z-shaped stroke (« per 2. huius », « per 20. primi Theodosij »), distinct
% from the plain 1; the 6 is often « b »-shaped; the 0 an « o ». Words
% written above the line with no deletion mark are given where they stand,
% after the word they sit over, with a \note{}. Figure letters are set in
% math ($abc$) on view 20; figures are described, not drawn. Tables are set
% as arrays in blocks of fifteen rows, figures as written, errors included,
% each block's inconsistencies listed in a \note{} (checked by adding the
% differences and against the sixtieths). Headings are set as \section* and
% \subsection*.
%
% Illegible: 10 (six of them the ends of verses hidden at the page edge on
% view 4). Uncertain readings: 24.
%
% Nothing here was checked against any printed edition of Maurolico or any
% other transcription.
\documentclass[11pt,a4paper]{article}
% The preamble shared by every transcription of the mathematicians' manuscripts in Gallica.
%
% It rests on one idea: **uncertainty compounds, it does not resolve.** A
% transcription that smooths over a doubtful word has destroyed the only thing
% it was made for — a reader must be able to tell, at every point, what was on
% the page and what was supplied. The apparatus macros below are therefore the
% critical apparatus, not decoration.
%
% The same commands are recognised by `scripts/render.mjs`, which produces the
% site's reading view, and by `scripts/tei.mjs`. A macro added here must be
% added there too, or rendering fails loudly — which is the wanted behaviour.
%
% Comments here are English, like the rest of the repository. The documents
% this preamble sets are French, and that is deliberate: see the skills.

\ProvidesPackage{germain}[2026/09/15 Transcriptions of mathematicians' manuscripts in Gallica]

\usepackage{amsmath,amssymb,amsthm}
\usepackage{mathrsfs}
% Kept from the parent preamble so the renderer's diagram support stays
% available; these manuscripts draw no commutative diagrams, and nothing here uses it.
\usepackage{tikz-cd}
\usepackage[english,french]{babel}
\usepackage{fontspec}

% --- Greek ------------------------------------------------------------------
% The text face stays fontspec's default, Latin Modern, which has no Greek:
% XeTeX drops every glyph it lacks with a line in the log and nothing on the
% page, and the Greek words the manuscripts quote vanished from the PDFs. So
% the Greek and Greek Extended blocks get a XeTeX character class of their own,
% and entering or leaving a run of it switches the family to CMU Serif — the
% same Computer Modern design, with polytonic Greek — and back. Only the
% family changes: a Greek word in \textit or \textbf keeps its shape and series.
% The transitions to and from every other class carry the tokens that class
% has towards an ordinary letter, so French spacing before « ; » or inside
% « » still holds after a Greek word. No grouping, so no brace or \endgroup
% can come out unbalanced; maths is left alone, since XeTeX inserts nothing
% there. Kept identical in `exercices.sty`.
\makeatletter
\newfontfamily\germainGreekfont{cmun}[
  Extension=.otf, NFSSFamily=germaingreek,
  UprightFont=*rm, ItalicFont=*ti, SlantedFont=*sl,
  BoldFont=*bx, BoldItalicFont=*bi, BoldSlantedFont=*bl]
\newXeTeXintercharclass\germain@greek
\def\germain@greekrange#1#2{%
  \@tempcnta=#1\relax
  \loop
    \XeTeXcharclass\@tempcnta=\germain@greek
    \ifnum\@tempcnta<#2\advance\@tempcnta\@ne\repeat}
\germain@greekrange{"0370}{"03FF}
\germain@greekrange{"1F00}{"1FFF}
\def\germain@greekprev{\rmdefault}
\def\germain@greekin{%
  \edef\germain@greekprev{\f@family}\fontfamily{germaingreek}\selectfont}
\def\germain@greekout{\fontfamily{\germain@greekprev}\selectfont}
\def\germain@greekjoin{%
  \XeTeXinterchartokenstate=\@ne
  \@tempcnta=\z@
  \loop
    \ifnum\@tempcnta=\germain@greek\else
      \XeTeXinterchartoks\germain@greek\@tempcnta=\expandafter{%
        \expandafter\germain@greekout\the\XeTeXinterchartoks\z@\@tempcnta}%
      \XeTeXinterchartoks\@tempcnta\germain@greek=\expandafter{%
        \the\XeTeXinterchartoks\@tempcnta\z@\germain@greekin}%
    \fi
    \ifnum\@tempcnta<4095\advance\@tempcnta\@ne\repeat}
% babel-french sets its own classes, and saves and restores the interchar
% state, each time a language is selected: join again after it does.
\addto\extrasfrench{\germain@greekjoin}
\addto\extrasenglish{\germain@greekjoin}
\AtBeginDocument{\germain@greekjoin}
\makeatother

\usepackage{geometry}
\geometry{a4paper,margin=25mm,marginparwidth=28mm}
\usepackage{marginnote}
\usepackage{xcolor}
\usepackage[normalem]{ulem}
\setlength{\emergencystretch}{3em}
\usepackage{hyperref}
\hypersetup{colorlinks=true,linkcolor=black,urlcolor=black,pdfborder={0 0 0}}

% --- Batch metadata ---------------------------------------------------------
% Taken as they stand by the rendering script, which shows them at the head of
% the reading view. They fix what the file claims to cover. The macro names are
% the parent project's, so that the scripts stay shared; \folder here means the
% volume's slug — `fr-9115` — and \pages counts Gallica views.

\newcommand{\folder}[1]{\gdef\grothFolder{#1}}
\newcommand{\batch}[1]{\gdef\grothBatch{#1}}
\newcommand{\pages}[2]{\gdef\grothFirst{#1}\gdef\grothLast{#2}}
\newcommand{\foldertitle}[1]{\gdef\grothFoldertitle{#1}}

% The holder's own shelfmark, printed as they write it: « Français 9115 ».
\newcommand{\shelfmark}[1]{\gdef\grothShelfmark{#1}}

% The dating, copied verbatim from the holder's catalogue — never inferred
% here. « XIXe siècle » is the BnF's; a pass that dates a leaf from its content
% says so in a \note{}, as its own inference.
\newcommand{\dating}[1]{\gdef\grothDating{#1}}

% The status notice, printed in the title block and repeated as a diagonal
% watermark on every page of the PDF. The manuscripts are in the public
% domain; what this line declares is not a right but a quality — an
% unverified first pass by a machine. The skills write the line; a file
% without it gets no watermark, so leaving it out is visible in the output.
\newcommand{\watermark}[1]{\gdef\grothWatermark{#1}}

\gdef\grothFolder{?}\gdef\grothBatch{}\gdef\grothFirst{?}\gdef\grothLast{?}%
\gdef\grothFoldertitle{}\gdef\grothShelfmark{}\gdef\grothDating{}\gdef\grothWatermark{}

% --- The résumé -------------------------------------------------------------
\newenvironment{resume}
  {\small\setlength{\parskip}{0.45em}}
  {\par\medskip\hrule\medskip}

% --- Keywords ---------------------------------------------------------------
% In English, closing the modernised reading's résumé: the modern vocabulary
% under which to search for what the leaves construct. The manifest extracts
% them to tag the volume — this line is the single source of the tags.
\newcommand{\keywords}[1]{\par\smallskip\noindent
  {\footnotesize\textsc{Keywords} --- \textit{#1}}\par}

% --- The critical apparatus -------------------------------------------------

% The Gallica view number, in the margin. This is the anchor that lets a
% disagreement over a formula be settled by looking at *one* image rather than
% a volume of 750. The site uses it to turn the facsimile as the transcription
% is scrolled. A view is one scanned image: usually one side of a leaf.
\newcommand{\page}[1]{%
  \marginnote{\footnotesize\color{gray!70}\textsf{v.\,#1}}%
  \label{page:#1}\ignorespaces}

% The BnF's pencilled foliation, where the leaf carries it and a pass can read
% it: « 198r ». This is what the literature cites — « ff. 198r–208v » — and
% the one line that turns a citation into a view. Recorded, never inferred:
% a folio a pass cannot read is not written.
\newcommand{\folio}[1]{%
  \marginnote{\footnotesize\color{gray!50}\textsf{f.\,#1}}\ignorespaces}

% The modernised reading's coarser counterpart to \page{}: which views a
% section is drawn from.
\newcommand{\pagerange}[2]{%
  \marginnote{\footnotesize\color{gray!70}\textsf{v.\,#1--#2}}%
  \ignorespaces}

% Illegible: never guessed.
\newcommand{\ill}{\textcolor{red!60!black}{[\,\ldots\,]}}

% A doubtful reading: the word is given, and flagged as uncertain.
\newcommand{\uncertain}[1]{\uline{#1}}

% An editorial addition — a missing letter, an implied word.
\newcommand{\add}[1]{\textcolor{blue!50!black}{[#1]}}

% Struck out by the writer. What was crossed out is part of the document.
\newcommand{\struck}[1]{\sout{#1}}

% The transcriber's note, on the state of the page or the mathematics.
\newcommand{\note}[1]{\par\smallskip{\small\color{gray!60!black}\textit{[#1]}}\par\smallskip}

% A marginal note of hers, distinct from the body of the text.
\newcommand{\marginal}[1]{\par\smallskip\noindent\hspace{1em}%
  {\small\color{gray!60!black}#1}\par\smallskip}

% --- Title ------------------------------------------------------------------
% The title block is French, like the documents themselves.

\AddToHook{shipout/background}{%
  \ifx\grothWatermark\empty\else
    \begin{tikzpicture}[remember picture, overlay]
      \node[rotate=52, scale=3.4, align=center, text=gray!18,
            font=\sffamily\bfseries]
        at (current page.center) {\grothWatermark};
    \end{tikzpicture}%
  \fi}

\AtBeginDocument{%
  \begin{center}
    {\large\bfseries Manuscrits de mathématiciens (Gallica) ---
      \ifx\grothShelfmark\empty\grothFolder\else\grothShelfmark\fi}\\[2pt]
    {\itshape\grothFoldertitle}\\[4pt]
    {\small vues \grothFirst--\grothLast
      \ifx\grothBatch\empty\else\ (lot \grothBatch)\fi}\\[2pt]
    \ifx\grothDating\empty\else
      {\small\color{gray!60!black}Datation du catalogue : \grothDating}\\[2pt]
    \fi
    \ifx\grothWatermark\empty\else
      {\small\bfseries\color{red!45!gray}\grothWatermark}\\[2pt]
    \fi
    {\footnotesize\color{gray!60!black}Transcription établie d'après les images IIIF de Gallica.
     Source gallica.bnf.fr / Bibliothèque nationale de France.\\
     Les lectures incertaines sont soulignées, les illisibles notées [\,\ldots\,],
     les ajouts entre crochets ; \textsf{v.} numérote les vues de Gallica,
     \textsf{f.} la foliotation au crayon de la BnF.}
  \end{center}
  \vspace{1em}}

\endinput


\folder{latin-17859}
\shelfmark{Latin 17859}
\batch{1}
\pages{1}{20}
\dating{1601-1700}
\watermark{Lecture automatique\\non vérifiée}
\foldertitle{Opuscules de François Maurolyco sur les mathématiques et l'astronomie.}

\begin{document}

\page{2}
\note{double page ; la page de gauche, dos du feuillet de garde, est blanche. Page de droite : la page de titre du recueil, numérotée à l'encre « 1 » en haut à droite. Dans la marge droite, d'une autre main et d'une plume plus grasse, l'ancienne cote « Nouv. acq. lat. 1099 », soulignée ; plus bas, le cachet rond « Bibliothèque impériale MSS ». Un petit trait vertical devant la première ligne.}

Theodosii Sphæricorum Elementorum libri tres. Ex traditione Maurolyci Messanensis Mathematici.

Menelai sphæricorum libri tres.

Autolyci de sphæra quæ movetur. Liber.

Theodosii de Habitationibus. Liber.

Euclidis Phoenomena brevissime demonstrata.

Demonstratio et praxis trium tabellarum scilicet sinus recti. Foecundæ, et Beneficæ ad Sphæralia triangula pertinentium.

Compendium Mathematicæ mira brevitate ex clarissimis autoribus.

Maurolyci de sphæra sermo.

Exscripta ex typis excusso codice Messanæ. In freto Siculo a Petro Spira mense Augusto Anno M.D.LVIII.
\note{cette ligne et les suivantes sont d'un module plus petit, en retrait à gauche.}

Bibliotheca \uncertain{I.U.} Thuanorum. mense Iunio. Anno M.DC.XXVIII. favore C.V.D.I.B. \uncertain{Haultini} mense decembri mutuo concesso Anno M.DC.XXVII.

\note{« I.U. » : deux lettres suivies d'un point chacune, la première en forme de J ; « Haultini » : le milieu du mot est surchargé. Au crayon, très pâle, sous ces lignes et d'une main moderne, en français : « Ce MS \uncertain{paraît} incomplet ». Le reste de la page est blanc.}

\page{3}
\note{double page ; page de gauche. Elle ne continue pas la page de titre : une lettre commence.}

\section*{D. Iohanni Cerda Methymnensium Duci ac Siciliæ Proregi Illustrissimo Maurolycus S.D.}

Tertio abhinc anno (Dux Illustrissime) hos sphæricorum libellos multis cum vigilijs emendatos D. CAROLO V. Imperatori maximo destinaram: neque ob temporum injuriam votum successerat. hoc interea mihi venit in mentem, ut intercessu ac favore tuo quod antea decreveram nunc perficerem. Quo enim patrono ad eam rem uti poteram nisi te quem Philippus eius filius Rex potentissimus, anno iam exacto nobis proregem dedit. Qui iam sicut totius Sicanicæ Reipublicæ atque militiæ curam magna cum laude sustines, ita litteratorum patrocinium suscepis. Vidimus te nuper inter equestres pedestresque copias armatum per Messenia littora barbaros hostes a Siculis finibus absterrentem: cum Turcorum numerosa classis Mamertino freto insidens Calabriam, igne ferroque popularetur. Esto nunc nobis apud alterum Augustum alter Mecoenas, et effice ut me serenus aspiciat. Qui quoniam Mathematicis operam dederam idcirco ab eo exordiri volvi quod maxime opportunum existimavi. Sphæralium quidem triangulorum scientia, quantum erat Geometricis et Astronomicis Studiis necessaria tantum in Gymnasijs hactenus neglecta iacuit. Nec erat ulterius tantus Academiæ speculationum tolerandus defectus. Itaque quoad potui conatus sum ut Theodosii et Menelai sphærica ex corruptissimis exemplaribus instaurata in lucem prodirent. His nonnullas lucubrationes meas et alia quædam ad idem negotium spectantia subiunxi, id tantum Prorex Excellentissime facias ut sicut me commendas ita quem sic suscipis libellum ad Serenissimum Cæsarem prospere dirigas: Ut sub tanto numine cætera quæ per indicem exponuntur coram exire non erubescant. Vale Messanæ Augusto mense M.D.LVIII.
\note{une tache d'encre entre « M.D. » et « LVIII »}

\subsection*{Index lucubrationum Maurolyci.}

Aliena.

Euclidis Elementa discussis errori\add{bu}s interpretum, cum additionibus quarumdam propositionum præsertim ad regularia solida spectantium.
\note{une tache d'encre couvre la fin de « erroribus »}

Theodosij sphærica Elementa lib. 3. Astronomicis principiis necessaria.

Menelai Sphærica lib. 3. multis demonstrationibus adaucta.

Apollonii Conica libri 4. et demonstrationibus et lineamentis instaurata.

Sereni Cylindrica libri duo.

Archimedis opera de dimensione circuli. De sphæra et Cylindro. De Isoperimetris. De momentis æqualibus De quadratura parabolæ.

\note{page de droite de la vue 3, numérotée à l'encre « 2 » en haut à droite ; elle continue la liste des œuvres d'Archimède (« De quadratura parabolæ » / « De sphæroidibus »).}
De sphæroidibus, et conoidibus figuris. De spiralibus in quibus multæ demonstrationes additæ et \uncertain{facilitas}.

Iordani Arithmetica. Rogerij Baconis et Iohannis Perspectiva cum adnotatibus quibusdam.
\note{« adnotatibus » porte un signe d'abréviation au-dessus du « t » : pour « adnotationibus ».}
Ptolemæi de speculis comburentibus libellus.

Autolyci de sphæra quæ movetur. Theodosij de habitationibus. Euclidis Phoenomena brevissime demonstrata. Problemata Aristotelis mechanica cum opportunis et notatu dignis additionibus.

Nostra vero sunt.

Prologi sive sermones quidam de divisione artium. de quantitate de proportione. De Mechanica autoribus. De sphæra de Cosmographia de Conicis. De solidis regularibus. De operibus Archimedis. De quadratura circuli. De Centris. De Instrumentis. De Calculo. De perspectiva. De Musica. de Divinatione. Arithmetica speculativa libris duobus. In quorum primo multa de formis numerorum a nemine hactenus animaduersa. In secundo autem Theoria et praxis rationalium et irrationalium magnitudinum per numerarios terminos abunde tractatur.

Arithmetica data libellis quatuor. Positionum et rei demonstrationes ad quatuor præcepta redacta. Sphæricorum libelli duo. In quibus multa a Menelao neglecta supplentur. Sphæra mobilis 8\textsuperscript{o} cap.\textsuperscript{us} Cosmographia de forma situ numeroque elementorum et coelorum olim Petro Bembo dedicata. Conicorum Apollonii quintus et sextus.

De Compaginatione solidorum regularium. Quæ figuræ locum impleant ubi Auerroïs ignorantia manifestissime constat.

De momentis æqualibus libri quatuor. In quorum postremo de centris solidorum ab Archimede omissis agitur. De quadrati Geometrici, Quadrantis, Astrolabii, speculatione fabricatioq; usuque.
\note{« fabricatioq; » : « q; » abrège « que » ; la fin de « fabricatione » n'est pas écrite}
De lineis Horarijs libelli tres In quibus tota huius modi linearum Theoria quoad situm colligantiam et descriptionem ipsarum plene tractatur. Nam lineæ horariæ a meridie coeptæ secant peripheriam quamdam in iis punctis in quibus eamdem tangunt lineæ horariæ ab occasu et ortu exortæ. Talis autem peripheria vel circulus est, vel ex conicis sectionibus aliqua scilicet parabole Ellipsis vel Hyperbole.

Photismi de lumine et umbra. Diaphana in tres libros divisa in quorum primo de perspicuis corporibus. In secundo de Iride. In 3\textsuperscript{o} de oculi structura et conspicillorum agitur. Quæstionum Arithmeticarum libri 3. Geometricarum libri 2. Astronomicarum libri tres. Annotationes \ill{}\note{un mot d'environ cinq lettres, à peu près « sioda », non lu.} et in diversas tabulas Canones. Compendium Mathematicæ ex præcipuis autoribus.
\note{le reste de la page est blanc.}

\page{4}
\note{double page ; page de gauche. Une seconde lettre dédicatoire commence ; la page de droite de la vue 3 s'achève sur l'index.}

\section*{AD C\add{arol}um V. Imperatorem maximum Francisci Maurolyci Siculi Abbatis Epistola.}
\note{le milieu du nom est effacé le long d'un pli du papier ; seuls « C » et « um » se lisent}

Cum diu multumque dubitassem Serenissime Cæsar Intempestivus ne te compellaturus accederem oblaturus simul Ingenii mei aliquem laborem, hoc præsertim tempore quo \struck{tantopere} magis ridiculus id tentarem quanto te major cura pacandis hostibus restituendaque Christianæ Reipublicæ intentum sollicitat, non habui tandem quo me converterem aut cui potissimum has Sphæricorum Theorias dicarem nisi tibi Cui summus orbis universi machinator terrestris sphæræ curam et regimen commisit. Iampridem enim hemisphærium nostrum partim tibi sponte militat, partim te aut dominum sentit aut bello fulminantem metuit. Reliquum vero quod antipodum est Argonautarum tuorum animosa industria quotidie pernavigat, non fabulosum sed verum tibi aureum vellus reportantium. Superest ut cuius sphæræ dominio potiris eius etiam exactam teneas notitiam. Adde quod huius nobilissimi corporis speculatio non erat nisi supremo principi dicanda tibi præsertim qui magnanimitate et munificentia, majores tuos longo superas intervallo, et qui ut animi celsitudinem ostenderes pariter et Philippum filium potentissimum faceres provincias et regna et anno præsenti Siciliam peractis Messanæ cerimoniis largissime donasti. Cumque sphæra cæteris figuris simplicitate æqualitate fortitudine capacitate formaque unitate præcellat, summoque quo potest representet archetypum, tibi etiam consideranda incumbit cui Iuppiter Imperium et aquilam communicat: mox inter sydera recepturus. Ratio ita postulat ut me de sphæricis loquentem non minus legas, quam regum ac principum tibi subditorum, ac pacem petentium litteras. Illi enim de urbe obsidenda aut provincia pacanda, nos de ambitu Coelorum ad quem pro Christo pugnans rebus prospere gestis aspiras agimus. Ad summam quanti pluris facienda est amplitudo Coeli, quam hic terra punctualis globus (in qua tamen nunc apostatas maiestate simul ac potentia domando nunc parcendo prostratis nunquam non labores) tanto pluris facere debes philosophorum de Coelestibus arcanis disserentium libellos. Animavit me ad istos olim suscipiendos labores Iohannes Vega tuus, quem aliquot ante annis Siciliæ præfeceras, qui et inter cætera quæ summa cum laude præstitit litterarum etiam curam et patrocinium non omisit. Contulit et favorem suum Simeon Vigintimillius Hieracensium Marchio litteratorum Mecoenas tunc Messanensium Strategus. una cum iuratis urbis patribus. Quorum suasu liber hic tantis itinerum spatiis peragratis verecundus tibi supplicaturus venit. Messanæ Iulio mense. M.D.LVI\ill{}.
\note{après « M.D.LVI », un ou deux jambages surchargés d'encre, non lus}

\subsection*{Eidem Augusto Dicatum carmen.}
\note{deux colonnes de distiques ; la fin des vers de la colonne de droite passe dans le pli de la reliure et ne se voit pas}

\begin{quote}
Arte Syracusia si \uncertain{nuclatum} \uncertain{viverat} orbem\\
Sydera perpetuo qui movet astra Deus.\\
Hoc monita exemplo Sicula vigilata per oras.\\
Ante tuos trepidat pagina Magne pedes.\\
Illius immensum numen reveretur Olympus:\\
Exterret trifido fulmine flamma reos.
\end{quote}

\begin{quote}
Te timet extanti qua cernit ab ethere Phoeb\ill{}\\
Quaque latet nautis cognita terra tu\ill{}\\
Inde pares aquilæ Clypeo geminantur a\ill{}\\
Dextra patris summi: leva ministra tu\ill{}\\
Illius est Coelum, sphæra hæc tibi nostra\ill{}\\
Divisum imperium cum Iove Cæsar ha\ill{}
\end{quote}

\note{page de droite de la vue 4, numérotée à l'encre « 3 » en haut à droite ; un texte nouveau y commence.}

\section*{Maurolyci Messanensis de Sphæra Sermo.}

Sicuti inter planas circulum ita inter solidas figuras sphæram maxima excellentia esse multis plane rationibus constat. Utraque enim in specie sua simplicitate similitudine partium, æqualite identitate loci fortitudine atque capacitate cæteris maxime præcellit.
\note{« æqualite » : ainsi écrit, sans signe d'abréviation visible, pour « æqualitate »}
Nam circulum unica linea, et sphæram unica claudit superficies. In circulo arcus similiter curvi, et in sphærâ portiones similiter convexæ. In utraque figura medium est ab extremis æqualiter remotum. Unde et utriusque latitudinem æquales diametri quaquaversum receptæ metiuntur. Utraque circa centrum circumvoluta eiusdem semper loci limitibus definitur. Unde utrique et motus facilitas et partium firmitas nullo obstante extrinseco maxima comparatur. Circulus demum inter planas isoperimetras et inter solidas figuras sphæra capacissima est. Quarum sex proprietatum quælibet satis esset ad figuræ diffinitionem: cum singula solis conveniant. Sed his accedit alia conditio mirabilis quod in neutra huiusmodi figurarum sit initium finemve reperire.

Quo fit ut non temere neque Sed consulto, ac summa dei providentia credamus factum ut mundus huiusmodi formam contraxerit, quo scilicet archetypum suum principio fineque carentem, qua posset similitudine imitaretur.
\note{« Sed » est écrit dans l'interligne au-dessus de « neque », qui n'est pas biffé ; les deux mots sont donnés dans l'ordre où ils se lisent.}
Sed et reliquæ dudum in ea forma proprietates memoratæ, universo maxime conveniunt. Simplicitas ut creatoris unitatem referret: similitudo ut partes totius formam participarent. Æqualitas ut æque gravia æqualiter a medio distarent. Identitas loci ne circumlatus vacuum relinqueret. sive ut totus in toto nunquam se non contineret. Fortitudo ut nihil resistentiæ, patiens quam facillime verteretur. Capacitas ut commodius omnia comprehenderet. Nec solum universus sed et partes orbes tam ethereæ quam elementaris regionis, et astra singula huiusmodi susceperunt formam. Sicut autem circulus a recta in plano altera extremitatum stante semel circumducta describitur: Ita et sphæra a semi\struck{diametro}circulo
\note{dans « æqualiter » et « partes », les finales « liter » et « tes » sont écrites dans l'interligne, au-dessus de la fin de la ligne ; « circulo » est écrit dans l'interligne au-dessus de « diametro » biffé. La phrase continue en tête de la vue 5.}

\page{5}
\note{double page ; page de gauche. Elle continue la page de droite de la vue 4 (« a semicirculo » / « super fixam diametrum »).}
super fixam diametrum semel revoluto designatur. Sed cum utriusque figuræ quoad diametros, chordas et angulos Euclides et Ptolemæus, quoad sphæralia triangula Theodosius et Menelaus, quoad areas et corpulentias Archimedes acutissimus proprietates abstrusas copiosissime pertractent. Nos ad ea descendemus quæ ad mundanorum orbium formam situm numerumque faciunt. Coelum enim esse sphæricum et motum eius circularem constat sensibili experimento, et quod stella circumvoluta semper eiusdem appareat magnitudinis. Nec non a facilitate et commoditate plurium motuum: et a cæteris huiusmodi figuræ dudum memoratis conditionibus. Neque aliis docemur rationibus astra singula esse sphærica. Sed neque lunaris crementi ac decrementi ratio constare posset si luminare ipsum alterius modi esset. Deinde quod globus hic, qui ex aqua et terra constat, sit rotundus patet ex hoc quod stellæ prius oriuntur orientalioribus aut prius ad meridianum perveniunt. Procedentibus autem sub uno meridiano constat ex variatione unius astri meridiana celsitudine arcusque diurni cremento decrementove. Cumque \ill{}\note{un mot de six lettres environ, à peu près « tetium », non lu.} ortivum intervalla et crementa sint interiectis terræ spatijs proportionalia, circularem rotunditatem esse necesse est. Præterea cum ex collibus locisque editioribus plus maris scopulosque quos e littore non cernimus prospectemus, quis dubitat convexum esse pelagus? Hoc idem indicat et terrestris umbra lunam in Ecclipsi percutientis rotunditas. Hoc æquilibrij corporum gravium æque ad centrum undique connitentium ratio requirit. Nam collibrata pelagi terræque facies cuiusmodi nisi æqualiter a centro remota et perinde sphærica esse potest? Ea enim collibrata iacent quæ sunt eiusdem altitudinis, et æqualiter alta ea quæ paribus ab universali centro spatijs distent. Ea demum a centro æque remota quæ in una sphærica superficie terminantur. Neque impedimento sunt montium eminentia aut valles: sicut nec asperitates parvæ ingentem lapideam pilam rotunditate privant. Cætera tria elementa liquiditate ac facilitate fluxus exactiorem consequuntur rotunditatem. Quod terræ

\note{page de droite de la vue 5, numérotée à l'encre « 4 » en haut à droite ; elle continue la page de gauche (« Quod terræ » / « tenacitas »).}
tenacitas non potuit, ut quæ inæqualis et cavernosa \uncertain{iam}\note{le mot se lit « ium », l'initiale en forme de J ; « iam » est offert.} in profundiora susceptis aquis, eminentias et insulas expeditas animalium terrestrium usui vitaque relinqueret: globumque totum maris ac lacuum supplemento adglomeraret, utque congeries ex utriusque natura corpore compaginata, repertis paribus collibratisque ponderibus idem gravitatis magnitudinisque centrum quod universale centrum est \uncertain{quieta} sortiretur. Cedat antiquorum ignorantia, quibus ingentes insulæ et immensæ regiones dudum repertæ fuerunt incognitæ. Cedat \uncertain{stultitia} diversas sphæras terræ marisque aut diversa centra ponentium.
\note{« quieta » : lettres surchargées. Le mot rendu « stultitia » se lit plutôt « Shelhitia » au trait ; la lecture est offerte par le parallèle avec « Cedat antiquorum ignorantia ».}
Deinde sicut aqua quæ levior est terra supernatat, sic altior circumfunditur aër, et aerem ignis levissimus complectitur. Sic graviora centro sunt viciniora: in centrum enim gravia tendunt, ut iam nobis ubicumque simus in terram calcantibus centrum illud infimum esse commune. Et quamvis antipodes erectos tamen utrimque stare mirandum aut incredibile nemini sit, nisi crasso ingenio laboranti, quique sit universale centrum minime intelligenti, Terram autem esse in centro ex his patet; quod stella eadem undevis eiusdem magnitudinis apparet. Quod semper dimidium Zodiaci sive coeli Semisphærij videmus. Quod umbrarum æquinoctialium termini in una recta linea sint, Quod nonnisi in oppositione luminarium fit Ecclipsis lunæ. Quod astro quoppiam oriente occidit oppositum. Ex his etiam concluditur Terram respectu firmamenti puncti vicem obtinere. Unde oculi observatoris a centro terræ distantia insensibiliter variat observationem instrumenti. Item cum perpendicula omnia sint axes horizontum et diametri terræ se invicem secent in centro universi. Iam terra cum universo idem centrum habet. præterea cum minima stellarum iampridem major quam terra firmamento collata punctum sit multo magis terra eodem respectu punctum erit. Superest nunc in hoc globo motum excludere. Non enim possibile est terram moveri motu recto sic enim relinqueret centrum quod contra sui naturam est. Quod moveretur motu circulari, tunc aut super alium axem ab axe mundi et sic variaretur altitudo poli unius loci, quod non fit. aut super ipsum axem mundi, et tunc nubes et volucres versus occidentem relinquerentur et lapis sursum iactus non recideret eodem. Quod cum minime fiat. Et motus sit gravi et quieto corpori contrarius, superest eam omnino esse immobilem. Motum autem coeli duplicem esse. Nam motus primus ab ortu ad occasum quo luminaria et astra oriuntur et occidunt, sentitur etiam a brutis. Motus autem huic contrarius ab occasu ad ortum constat ex vario solis ortu et ex varia velocitate Lunæ, ac planetarum versus ortum delatorum. Quoniam igitur coelum primum, luminaria, et quinque erratici totidem diversis

\page{6}
\note{double page ; page de gauche. Elle continue la page de droite de la vue 5 (« totidem diversis » / « motibus aguntur »).}
motibus aguntur. Ideo coelos pauciores esse quam octo nequaquam posse. Singuli enim motus singulos sibi postulant orbes. Deinde cum perceptum esset a Ptolemæo stellas fixas præter diurnum alio tardissimo motu versus Orientem et mox a Thebitis motu trepidationis ferri, Oportuit nonum coelum addere. Cum vero Alphonsus trepidationi motum longitudinis addidisset sic opus fuit Coelo decimo nisi quis duos hos quasi deferentes octavi faciat. Dispositio autem et ordo sphærarum sumitur a quantite secundarij motus.
\note{« quantite » porte un signe d'abréviation au-dessus du « t » : pour « quantitate ».}
Nam tardiori motui congruit superior locus Unde octava sphæra suprema est Lunaris infima. Item ex Ecclipsibus ut quoniam Sol obtegitur nobis a luna, idcirco ea superior est. Deinde ex diversitate aspectus veri visualisque quæ major existens astrum arguit inferius et maxima lunam infimam. Adhuc ex convenientia. Congruum enim fuit Solem inter planetas principem locari medium martemque ac venerem soli collaterales poni: et qui deferentum auges una cum solis auge sortiti et epicyclorum magnitudine cæteris præstantes majori tractu solis viciniam combustionisque iacturam vitare possent. Nam quod Saturnus Iove, venus Mercurio sit superior, nulli dubium. Præterea ne ingens intervallum inter lunarem ac solarem sphæram superfluum dimitteretur: dandum fuit veneris ac Mercurij coelis. Talis est coelestium sphærarum ordo a vetustissimis Astronomis positus, et recentiorum iudicio comprobatus.

Circulus itaque in sphærâ maximus super cuius axe polisque fit motus primus æquinoctialis est inter parallelos, eo motu descriptos. Hunc oblique secat Zodiacus Solis et planetarum orbita, qui tangit binos Æquatoris parallelos æquales qui Tropici dicuntur limites secundi motus. Coluri per polos mundi transeunt: unus per polos etiam Zodiaci et contactus tropicorum, reliquus per æquinoctialia puncta. Horizon est circulus manifestum hæmisphærium ab occulto disterminans cuius polus vertex loci vel zenit dicitur. Rectus horizon per mundi polos. Obliquus præter polos transit. Meridianus per mundi et horizontis polos incedit. verticalis per horizontis et meridiani polos Circulus latitudinis per Zodiaci polos. Circulus declinationis per æquatoris polos, uterque per astri locum ducitur. In illo astri latitudo, in hoc declinatio computatur. Circulus altitudinis per polos \struck{altitudinis} horizontis, in quo astri altitudo super horizontem mensuratur. Latitudo loci est arcus meridiani inter Zenit et æquatorem, sive altitudo poli. Longitudo vero loci in æquatore numeratur a meridiano extremi Occidentis usque ad meridianum ipsius loci. Describentes enim mundum

\note{page de droite de la vue 6, numérotée à l'encre « 5 » en haut à droite ; elle continue la page de gauche (« Describentes enim mundum » / « septentrionalia »).}
septentrionalia ponimus superiora, et ordinem scribendi et successionem signorum sequamur. Horizon itaque rectus quoniam per polos transit, omnes parallelos per motum primum descriptos singulos per æqualia secat, Unde ibi omnes stellæ oriuntur et occidunt, et æquinoctium est perpetuum, Obliquus autem horizon tangit duos æquatoris parallelos invicem æquales Inter quorum extantem stellæ nunquam occiduntur. Inter occultum nunquam oriuntur Et obliquior horizon tales parallelos majores et tales stellas plures habet. nam reliquæ stellæ, ortum patiuntur et occasum Sed eorum paralleli secantur inæqualiter ab horizonte. Arcusque diurni versus polum manifestum majores nocturnis, versus reliquum minores, et eo quo remotiores ab æquatore. Coalterni autem arcus æquales. Et obliquior horizon majorem patitur inæqualitatem. Habitantibus ergo intra tropicos Sol in anno transit bis per zenit, et zodiacus bis in die fit rectus horizonti. Habitantibus vero sub tropico utrumque dictorum accidit Cæteris nunquam. Quare qui intra Tropicos amphiscij dicuntur quod illis umbra meridiana utroque declinet. Cæteri heteroscij quod in partem alterutram. Maximus autem dies, et maxima nox cum fiat apud Tropicorum contactus, Iam horizon habens polos in arctico et antarctico parallelis, qui per Zodiaci polos describuntur tanget tropicos et ideo tum maximus dies, quam maxima nox erit inter circulus scilicet 24. horarum. Et \uncertain{instans} pro nocte vel die.
\note{« instans » : le mot se lit à peu près « insterps » ; la lecture est offerte par « et instans pro nocte » de la vue 12. « nocte » est repassé à l'encre.}
Horizon vero polos habens inter arcticum vel antarcticum circulum et mundi polum tanget duos parallelos Tropicis majores qui abscindent hinc et inde de zodiaco duos arcus æquales, quorum unus nunquam occidet, reliquus nunquam orietur. Quamdiu igitur Sol erit in illo fiet continuus dies, quamdiu vero in hoc nox continua: tot scilicet dierum quot gradu\add{s} aut tot mensium, quot signa amplectitur arcus abscissus.
\note{une tache d'encre sur « quot » et sur la fin de « gradus »}
Nam horizon polum habens in polo mundi est et æquator habens semicirculum Zodiaci extantem et semicirculum latentem. Et idcirco dies continuus semestris, et nox itidem. Zodiacus autem circulus, quoniam in motu diurno, convertitur super alienos polos. Ideo eius partes non æqualiter descendunt et ascendunt, sicut partes æquatoris. In horizonte tamen recto quatuor arcus zodiaci æque remoti a punctis æquinoctialibus oriuntur et occidunt æquis temporibus. Arcus vero remotior in majori. In horizonte autem obliquo hanc collationem facies in semicirculo zodiaci ascendente. Ascensio igitur signi vel arcus Zodiaci est arcus æquatoris cum eo simul oriens. Descensio

\page{7}
\note{double page ; page de gauche. Elle continue la page de droite de la vue 6 (« Descensio » / « simul occidens »).}
simul occidens. Recta quidem in recto, Obliqua in obliquo horizonte. Et oppositæ peripheriæ ex diametro oriuntur et occidunt. Et cuiusvis peripheriæ ascensio et descensio coniuncta ubique conficiunt eumdem arcum. Item arcus Zodiaci cuius limites æqualiter distant a solstitio eamdem ubique habent ascensionem sive descensionem. Arcus vero Zodiaci existens in semicirculo zodiaci ascendente vel majorem ibi portionem habens minorem sortitur ascensionem descensionem vero majorem. Præterea quoniam tum noctu quam interdiu exoritur et occidit ubique semicirculus zodiaci Iam in recto horizonte quivis duo semicirculi zodiaci temporibus æqualibus oriuntur aut occident. In reliquis locis diversificabuntur tempora pro quantitate arcuum diurnorum et nocturnorum. Unde habentibus zenith in arctico vel antarctico continget semicirculum zodiaci ascendentem in instanti exoriri et in spatio 24 horarum occidere. Et in reliquo fieri contrarium. Ortus autem et occasus astrorum considerantur etiam respectu solis. Oriente namque sole stella oriens fuit ortum, occidens autem occasum matutinum.
\note{« fuit » est écrit ; le sens demande « facit »}
Occidente vero sole utrumque vespertinum. Unde ortus matutinus et occasus vespertinus possunt contingere eodem die. Et similiter ortus vespertinus, et occasus matutinus si astrum sit in zodiaco. Stellæ tardiores sole desinunt videri vesperi propter accessum solis ad eas, et dicitur postrema fulsio vespertina. Postea præterito a sole incipiunt apparere mane ante solem orientem, et dicitur prima fulsio matutina. Contrarium facit luna quia velocior est sole. Ipsa enim accedens ad solem mane videri desinit et dicitur postrema fulsio matutina. Et paulo post præteriens solem incipit apparere vesperi et dicitur postrema fulsio vespertina, quod eodem die visa fuit quandoque fecisse.
\note{« postrema fulsio vespertina » est écrit là où le parallèle avec ce qui précède attend « prima » ; « vespertina » est repassé à l'encre}
Venus autem et Mercurius qui quandoque velociores sunt sole quandoque tardiores, et retrogradi, utramque passionem geminant, Nam circa augem Epicycli faciunt id quod luna et in opposito augis id quod astra superiora, Et sciendum quod stellæ ad manifestum polum magis accedentes minus occultantur a sole: nam propter magnitudinem arcus eorum diurni mane possunt prævenire solem Orientem et Occidentem subsequi, sic ante postremam visionem vespertinam incipiunt fulgere matutinæ, Signa physica subtenduntur a sex hexagoni lateribus et singula bifariam secta faciunt duodecim signa communia, et hoc singula in 30 gradus idcirco et annus in 12 menses. et mensis in 30 dies. Ut sicut circulus anno ita signum mensi, et gradus die responderet. Congruit huc et lunaris motus: nam sicut

\note{page de droite de la vue 7, numérotée à l'encre « 6 » en haut à droite ; elle continue la page de gauche (« nam sicut » / « lunatio »).}
lunatio mensem fere metitur, sic quadratura et plenilunium quadrantes, hoc est hebdomadas distinguant, quamvis non ad amussim: Ut totus Kalendarum computus pendeat a luminaribus. Dies est spatium temporis quo integre convertitur æquator addito arcu motus Solis interim peracto respondente Unde propter inæqualitatem talis arcus et motus solaris erunt dies inter se vulg. Inæquales, qui et vulgares vel apparentes vocantur. Cumque in toto anno additamenta additamenta diurna collecta perficiant circulum, Iam circulus æqualiter rursum divisus in numerum dierum anni exhibebit additamentum mediocre unius diei, sic dies reducuntur ad æqualitatem, qui dies æquales et astronomici dicuntur et congrui calculo motuum.
\note{« vulg. » en fin de ligne, puis « Inæquales » à la ligne ; « additamenta » est écrit deux fois de suite}
Vulgaris autem et Astronomici dierum differentia est æquatio dierum, quæ in dies collecta sensibilis redditur a tempore vulgari propterea subtrahenda. hora æquinoctialis est 24. pars diei. Hora vero naturalis vel temporalis, est 12 pars diurna diei, et nocturna noctis, per quas distribuantur dominia Planetarum secundum Babylonios et Hermetem. Perioeci sunt incolæ eiusdem paralleli, quibus eadem latitudo et arcuum loci dispõo.
\note{« dispõo » : abrégé, pour « dispositio »}
Antoeci sunt incolæ oppositorum et æqualium parallelorum quibus æqualis latitudo et arcuum, sed pro oppositis polis locisque dispositio, Antichthones sive Antipodes sunt qui non solum antoeci sed ex diametro \uncertain{stantes} oppositi, et qui eundem horizontem et diversa hæmisphæria sortiuntur: Unde quicquid apud hos apparet, aut oritur apud illos delitescit et occidit et e contrario, Exordium diei sumitur a meridiano ad evitandam multiplicem horizontum varietatem: in calculo et dierum æquatione. quandoquidem meridianus meridianum eodem temporis intervallo præcedit. Non sic autem horizon horizontem, nisi tantum in perioecis. Climatum distinctio a Ptolemæo fit secundum crementum maximi diei per horæ dimidium: parallelorum vero per quadrantem. Linea veri loci stellæ a centro mundi per centrum stellæ ducitur ad concavam \struck{coeli} primi mobilis superficiem, in qua per talem lineam circulos omnes
\note{les deux dernières lignes, depuis « stella ducitur », sont d'une encre plus pâle}

\page{8}
\note{double page ; page de gauche. Elle continue la page de droite de la vue 7 (« circulos omnes » / « descriptos »).}
descriptos intelligimus, Inæqualitas motus in planetis salvatur per eccentricos et etiam per Epicyclos, per quos salvantur salvantur Stationes regressus et latitudines.
\note{« saluantur » est écrit deux fois de suite}
Planetæ corpus fertur in Epicyclo in superiori parte secundum ordinem signorum in inferiori contra. Luna vero in contrarium, Epicyclus vehitur ab Eccentrico versus orientem. Sed soli solis est Eccentricus sine Epicyclo, sive concentricus cum Epicyclo. Recta per mundi ac deferentis centrum procedens indicat punctum maxime remotum a centro mundi in peripheria deferentis, et maxime propinquum: scilicet augem, et eius oppositum. Motus autem æqualis in Sole sive medius fit respectu deferentis: in deferente lunæ respectu centri mundi. in deferentibus reliquorum respectu centri æquantis quod in tribus superioribus et venere tanto superius est centro deferentis, quanto hoc centrum superius centro mundi. In Mercurio autem medium est inter centrum mundi centrumque parvi circuli, hoc est in ipsa circuli huius peripheria in qua contra signorum successionem, fertur centrum deferentis ea velocitate qua deferens in contrarium movetur. Sic centrum lunaris deferentis fertur in peripheria parvi circuli mundo concentrici contra successionem quantum Epicyclus in contrarium. Ab infimo autem dicti circuli puncto linea ducta per centrum Epicycli terminat augem eius mediam. In cæteris autem a centro æquantis. Nam augis vero index a centro mundi educitur. Argumentum medium in Epicyclo ab auge media, verum a vera usque ad planetam numeratur. Linea vero a centro Eccentrici in sole, a centro mundi in Luna, a centro æquantis in cæteris, per centrum solis vel Epicycli ducta, æquidistantem: ibi habet rectam a centro mundi medii motus determinatricem, et in lunâ est ipsamet cum a centro mundi exeat, Argumentum in sole Centrum in reliquis est distantia ab auge deferentis, secundum motus ordinem numerata, Æquatio sive in Epicyclo sive in deferente est motuum medii et veri differentia, quæ apud augem vel oppositum eius nulla est maxima vero in mediis sive in contactuum locis. Auges deferentium et in Mercurio aux æquantis moventur motu octavæ sphæræ Nam aux defferentis Mercurij ultra motum octavæ accedit, et recedit intra lineas a centro mundi per contactus parvi circuli ductas. Aux vero eccentrici lunaris ita contra successionem signorum movetur ut linea medii motus solis sit media inter augem et epicyclum, vel simul cum eis ut in coniunctione et oppositione. Unde in quadraturis centrum Epicycli semper est in opposito augis. Tres autem superiores tantum absunt

\note{page de droite de la vue 8, numérotée à l'encre « 7 » en haut à droite ; elle continue la page de gauche (« tantum absunt » / « ab auge Epicycli »).}
ab auge Epicycli quantum sol ab Epicyclo. Unde in auge sunt cum sole. Et in augis opposito sunt et soli oppositi. Dum vero inferiores communem habent medium cum sole motum. Capientur autem autem æquationes per centrum vel argumentum.
\note{« autem » en fin de ligne, répété en tête de la suivante}
Æquatio argumenti pro luna apud augem deferentis, pro cæteris apud longitudinem mediam notatur in tabulis. pro aliis autem locis: per minuta proportionalia, et diversitatem diametri, in luna simplicem, in cæteris duplicem rectificatur. Est autem diversitas diametri æquationum correlativarum differentia. Demum numeratis mediis motibus per adiectionem, vel subtractionem æquationum veri eliciuntur, Statio prima in Epicyclo est punctum in quo \struck{remisso motu directo, redit ad directionem} planeta urgente Epicycli motu incipit regredi. Statio secunda punctum in quo remisso motu dicto redit ad directionem. quæ puncta per lineas a centro mundi Epicyclum secantes determinantur. Diversitas aspectus astri est veri visique loci differentia, in circulo altitudinis computata. Hinc diversitas aspectus in longitudine et latitudine notescit. Luna in oppositione Solis interiecta terra Sol vero \struck{inter} interpositione lunæ sibi coniuncta defectum luminis patitur, quanquam sol tegitur non deficit. Quod fit in nodis vel iuxta quæ sunt puncta sectionum ecclipticæ cum lunari deferente. Circa quos termini assignantur intra quos Ecclipsis est possibilis, et rursum termini arctiores, intra quos Ecclipsis est necessaria, Cum vero Ecclipsis Solis fit coniunctio visibilis in ea totus calculus fit per visos locos, arcusque qui ex diversitate aspectus notescunt. Est autem Ecclipsis totalis aut partialis sine mora vel cum mora. Nodi autem lunæ moventur super axe Zodiaci contra successionem. Unde Ecclipses variant loca. Planum Epicycli lunæ in plano deferentis iacet Unde luna unicam \struck{habet} latitudinem ex deviatione deferentis patitur: quæ per distantiam a nodo quod est argumentum latitudinis captatur. In quolibet autem trium superiorum deferens declinat ab Eccliptica et aux deferentis septentrionalis In marte quidem apud punctum maximæ declinationis: in Iove procedit per 20 gradus In saturno sequitur per 50. Dumque centrum Epicycli sistitur in nodis eius planum iacet in plano Ecclipticæ: et diameter Epicycli per augem veram in communi sectione planorum, Inde discedens Epicyclus inclinat \ill{} dictum diametrum, ita ut aux \struck{Ecclipticæ} Epicycli planis Ecclipticæ et deferentis interiaceat, et diameter secunda longitudinum mediarum iugiter æquidistet ecclipticæ, et maxima inclinatio fit in puncto maximæ declinationis et inde rursum decrescit usque ad alium nodum.
\note{le mot non lu avant « dictum » se présente à peu près comme « sontim »}
Sic planetæ latitudo ex inclinationibus deferentis et Epicycli componitur: In duabus vero inferioribus deferens deviat ab Eccliptica ipsamque præteriens accedit et recedit super nodos tanquam polos et aux æquantis in puncto maxime declinante. Epicyclus quoque inclinat primam diametrum et reflectit secundam, Deviatio

\page{9}
\note{double page ; page de gauche. Elle continue la page de droite de la vue 8 (« Deviatio » / « autem deferentis »).}
autem deferentis et reflexio Epicycli nullæ sunt, centro Epicycli in nodis, maximæ vero in punctis maxime deviantibus; constituto. Contra inclinatio Epicycli hic nulla, ibi maxima. In locis medijs crescunt et decrescunt, hac lege ut centrum Epicycli veneris semper sit in semicirculo deferentis septentrionali, Mercurii vero australi. Ita ut dum centrum Epicycli recedit ab auge deferentis. aux epicycli inclinetur a deferente in venere, ad septemtrionem. In Mercurio ad Austrum: post autem oppositum augis deferentis fiat contrarium. Item ut dum \struck{centrum Epicycli recedit ab auge deferentis, aux epicycli inclinetur a deferente in venere ad Septemtrionem. In Mercurio ad austrum. Post} Epicyclus recedit a nodo anteriori versus augem, semidiameter orientalis Epicycli reflectatur in venere ad Septemtrionem, In Mercurio ad austrum post autem reliquum nodum accidat huius contrarium. Sic tum veneris quam Mercurii latitudo ex deviatione deferentis inclinatione et reflexione Epicycli componitur. Latitudines autem in tabulis notatæ sunt maximæ quæ capientur per centra et argumenta et rectificantur ad numerum minutorum proportionalium, per distantiam Epicycli a loco maximæ latitudinis sumptorum. Ex hos rursum probabiliter arguitur venus Mercurio superior quandoquidem septentrionalitas australitatem et orientalitas occidentalitatem dignitate præcedit.

Octava vero sphæra præter motum diurnum, movetur motu duplici, scilicet trepidationis quo sibi proprius est et longitudinis quem habet virtute nonæ sphæræ. Nam diameter coniungens initia arietis et libræ Octavæ sphæræ ita circumducitur ut medio puncto quod est in mundi centro fixo manente extrema ferantur in peripheriis parvorum circulorum quorum poli sunt initia arietis et Libræ nonæ sphæræ. Hac lege ut sectio duorum talium zodiacorum semper fiat in initijs Cancri et Capricorni nonæ sphæræ quæ revolutio fiat in spatio 7000 annorum. Motus autem longitudinis fit virtute nonæ sphæræ secundum ordinem signorum cuius revolutio completur in spatio 49000 annorum. Ex his duobus motibus componitur motus octavæ sphæræ in stellis fixis. Et motus trepidationis numeratur in peripheria parvi circuli a puncto maxime septemtrionali et per ipsum quasi argumentum, capitur æquatio quæ est ad summum 9 graduum, quanta parvi circuli semidiameter, addenda sive subtrahenda motui longitudinis ad rectificandos augium deferentum, et fixarum motus, iuxta commentum Alphonsi, quem Astronomorum communis Academia sequitur. Ex prædictis apparet primatus et dignitas Solis inter astra quippe qui non solum regulam motus superioribus et inferioribus planetis præscribit, sed etiam diurno motu gradus menstruo menses, et annuo communem motuum describit orbitam, cujus respectu tam longitudines quam latitudines erratium et fixarum stellarum computantur. Asterismi autem sive constellationes a majoribus nostris instituta sunt vel a forma

\note{page de droite de la vue 9, numérotée à l'encre « 8 » en haut à droite ; elle continue la page de gauche (« vel a forma » / « ut sagitta »).}
ut sagitta, corona, fluvius serpens. vel ab influxu, ut ara quoniam facit sacrificos, Delphin natatores. vel a natura seu dominio ut scorpius, quoniam scorpionum terrestrium naturam sapit. seu quod ipsis dominetur. Pisces quod signum aquaticum, vel ab historia seu memoria rerum gestarum ut Hercules, Perseus Ophiuchus. vel a placitis Principum ut crinis Berenices Reginæ a Conone Philosopho et Callimacho poeta in cauda Leonis locatus. Denique digestis in aliquot simulacra stellis, iam singula per partes et membra distributa quasi per localem memoriam facilius retenta promptius proferrentur. Differunt tamen autores in quibusdam ut signum quod apud Higinum Cygnus in asterismis Ptolemæi et Alphonsi Gallina, et avis dicitur. Quod ab illis Aquila ab \struck{illis} his vultur volans: Quæ ab illis lyra ab his vultur cadens appellatur: quamvis circa naturas constellationum non discrepent.
\note{« his » est écrit dans l'interligne au-dessus de « illis » biffé}
Divisio duodecim domorum secundum Iohannem de Monteregio fit in æquinoctiali, secundum Campanum in circulo verticali. Sed melius esset hanc divisionem et omnem directionum et profectionum calculum referre ad orbitam solis ut Ali, Abraham, Dorotheus docent, et recentioribus Trapezuntio et Cardano placuit: et Schonerus multis argumentis comprobavit. Quando duo mobilia eodem versus feruntur, superatio quando in diversas adgregatum motuum ad distantiam mobilium est sicut tempus superationis vel adgregati ad tempus quousque mobilia coeant numeratum. Quando duo mobilia in periferia circuli feruntur inæqualiter, tunc si eodem versus differentia si in diversas adgregatum terminorum erit numerus punctorum circulum per æquos arcus dividentium, in quibus punctis mobilia semper coeunt. Terminos intelligo minimos numeros motuum proportionem representantes. Si vero motus fuerint incommensurabiles, puncta coniunctionum in infinitum variabuntur. Quare si motus Coelestes sint tales falluntur Platonici, eamdem astrorum dispositionem ad tempus quoddam redire putantes. Zodiaci obliquitas ex altitudinum solis meridianarum collatione. Solis ingressus in æquinoctium ex altitudine ipsius meridiana. Motus eius quantitas ex collatione duarum observationum, Solis eccentricitas et locus augis ex observatione trium locorum. Locus lunæ per Lunares Ecclipses Locus eius in Epicyclo per motus eius varietatem. Intervallum duarum Lunarium Ecclipsium, in quibus luna ab eodem nodo et versus eamdem partem, æqualiter distat in eodem Epicycli loco sit complectitur integras argumenti latitudinis reditiones, et integros menses Lunares. Hinc latitudinis et capitis draconis motus quantitas. Diameter Epicycli per tres lunæ locos in tribus Ecclipsibus consideratos. Ex variatione diversitatum propter Epicyclum conicitur diversitas propter eccentricitatem,

\page{10}
\note{double page ; page de gauche. Elle continue la page de droite de la vue 9 (« propter eccentricitatem » / « et ipsa eccentricitas »).}
et ipsa eccentricitas. Ex tertia variatione aux Epicycli media. Ex differentia loci visi et veri eliciuntur lunæ a terrâ remotio respectu diametri terræ. Hinc ex ecclipsi solis et ex visualibus diametris luminarium, et umbra comprehenduntur reliquæ distantiæ, et corporum veræ diametri et proportio. Locus astri vel planetæ, per instrumentum vel collationem ad fixas noscitur. Intervallum duarum observationum, in quarum utraque planeta in eodem Zodiaci loco, et in eadem a sole distantia fit, continet integras deferentis et Epicycli revolutiones. Hinc motuum quantitas. Locus augis inferiorum per eius æquales et contrarias distantias maximas a sole. Unde diameter Epicycli. In superioribus autem locus Epicycli est idem cum loco planetæ in oppositione Solis. Et ex tribus locis Epicycli colligitur eccentricitas, locus augis et epicycli diameter Centrum æquantis et eccentricitas in illis per adgregatum duarum elongationum maximarum, et contrariarum a sole. In his autem supponitur idem centrum deferentis medium inter centra mundi et æquantis, et positioni consonat experientia Si nodi lunæ sint in punctis æquinoctionalibus, ex altitudine eius meridiana elicitur deviatio deferentis. Ex maximis vero latitudinibus planetarum, et ex distantijs et diametris Epicyclorum patent inclinationum anguli, Motus instauratis observationibus rectificantur, et \uncertain{radix} statuitur.

FINIS.

\note{« radix » : le mot se lit à peu près « ruedix », sans certitude.}

\section*{THEODOSII SPHÆRICORUM ELEMENTORUM EX TRADITIONE MAUROLYCI LIBER PRIMUS. II\textsuperscript{us} et III\textsuperscript{us}}
\note{« us » rend un signe en forme de 9 placé en exposant après les chiffres romains}

DIFFINITIONES.

Sed hi collati sunt cum veteri editione Veneta anni 1518.
\note{cette ligne est d'une écriture courante plus petite que les titres en capitales qui l'encadrent. Ni les définitions ni le texte de Théodose ne suivent : le bas de la page est blanc.}

MENELAI SPHÆRICA EX EIUSDEM TRADITIONE

Alibi descripta.
\note{« Alibi descripta » : le texte de Ménélas n'est pas copié ici. Le reste de la page est blanc.}

\note{page de droite de la vue 10, numérotée à l'encre « 9 » en haut à droite. Elle ne continue pas la page de gauche : elle s'ouvre sur la fin d'une proposition dont le début n'est pas dans ce lot. Le sens la rattache à la page de gauche de la vue 48 (lot 3), « Autolyci de sphæra quæ movetur ex traditione Maurolyci », dont la proposition 4 s'arrête sur « et per primam huius puncta singula » : la phrase se lit « puncta singula suos seorsum parallelos… describent », et les propositions 5 à 9 qui suivent sont celles du même livre sur la sphère en mouvement. C'est une reconstitution de cette transcription ; le titre de la vue 2 place de même Autolycus après Ménélas. Les numéros « Propositio 5 » à « 9 » sont en marge à gauche ; les démonstrations sont d'un module plus petit.}

suos seorsum parallelos, semper in alterutro hæmisphæriorum describent.

\textbf{Propositio 5.} Si per polos sphæræ circulus manens diffiniat apparens, et non apparens cuncta in superficie sphæræ puncta ipsa evoluta et occidunt et oriuntur, et æquali tempore morantur super horizontem, et sub horizonte.

Nam talis circulus manens per 20. primi Sphæricorum Theodosij secat per æqualia singulos parallelos hoc est in semicirculos.\note{« 20. » : le 2 de cette main, dans les renvois, est un petit trait en forme de z (même tracé dans « per 2. huius » et « per 6. 2. ») ; il est distinct du 1, un simple jambage.} Quare per 2. huius per æquale tempus punctorum unumquodque feretur utrimque a circulo secante.

\textbf{Propositio 6.} Si in sphæra maior circulus manens definiat quod apparens est sphæræ, et quod non apparens obliquus existens ad axem attinget binos circulos æquales, et parallelos invicem, et eorum unus ad apparentem polum semper erit apparens, alter autem ad latentem semper latens.

Constat hoc manifeste ex 10. Secundi Sphæricorum Theodosij.

\textbf{Propositio 7.} Si circulus in sphæra fixus apparens ab occulto distinguat obliquus existens ad axem circuli qui ad angulos rectos axi, in eisdem semper punctis horizontis ortus et occasus faciunt, et similiter inclinantur ad horizontem.

Nam cum circulus fixus constet semper in eodem loco, et circuli super axe suo versati semper in suo singuli plano iaceant, fit ut neque peripheriæ locum aliquando commutent: et perinde secent fixum, in iisdem semper punctis. Inclinatio quoque eorum una est, sequitur enim inclinationem axis communis.

\textbf{Propositio 8.} Si circulus major in sphæra in sphæra fixus apparens ab occulto dirimat inclinatus ad axem, quicumque circulus major contingit duos circulos parallelos æquales semper videlicet, apparentem semperque occultum, quos horizon contingit, evoluta sphæra congruit horizonti.
\note{« in sphara » est écrit deux fois}

Nam dum versatur sphæra puncta contactuum feruntur semper in peripherijs dictorum parallelorum, et perinde contactus dicti circuli majoris coniunguntur contactibus horizontis, et circulus ipse counitur horizonti.

\textbf{Propositio 9.} Si in sphæra major circulus obliquus ad axem definiat manifestum ab occulto, quod ex simul orientibus punctis est polo apparenti propinquius, posterius occidit, Quod autem ex simul occidentibus dicto polo vicinius prius oritur.

\page{11}
\note{double page ; page de gauche. Elle continue la page de droite de la vue 10 (« vicinius prius oritur » / « Nam punctum »).}
Nam punctum polo manifesto vicinius habet per 24. secundi Theodosij, majorem arcum super horizontem, et perinde si simul oritur cum puncto remotiori a dicto polo posterius occidet, et si simul occidat, iam prius exortum est per 2. huius.

\textbf{Propositio. 10.} Si in sphæra maximus orbis obliquus ad axem definiat apparens sphæræ, et latens, circulus qui per polos sphæræ bis rectus fit horizonti in uno sphæræ ambitu.

Patet quoniam talis circulus bis transit in uno ambitu per polos horizontis: quare per 20. primi Sphæricorum Theodosij bis eum orthogonaliter secabit.

\textbf{Propositio. 11.} Si in sphærâ major circulus obliquus ad axem definiat apparens sphæræ et latens, alius vero major circulus parallelos majores attingat, aut quos horizon semper apparentem semperque occultum tangit, per omnem horizontis peripheriam parallelis quos attingit interpositam ortus et occasus \struck{fiet} facit.

Patet quoniam omnia puncta talibus parallelis interiecta oriuntur et occidunt, apud peripherias horizontis iisdem interiacentes, Quare et tota talis circuli majoris peripheria id idem facit.

\textbf{Propositio. 12.} Si in sphærâ manens circulus delatum aliquem eorum, qui in sphærâ semper per æqualia secet neuter autem ipsorum, ad rectos fuerit angulos axi, neque per polos sphæræ, uterque ipsorum erit major.

Nam si uterque sit circulus \struck{m} maior: manens non potest semper bifariam secare delatum nisi manens ad rectos sit axi. Si manens sit major ac delatus minor, non potest semper bifariam secare delatum, nisi existentem ad ad rectos axi, quod est contra hypothesim. Si manens sit minor ac delatus major. hoc esset contra 17. primi Sphæricorum Theodosii. Superest ut omnino sint ambo majores.
\note{« maior » est écrit après un début de mot biffé (« m » suivi d'un jambage) ; le raisonnement demande ici « minor ». « ad » est répété à la ligne.}

\section*{Theodosii De Habitationibus ex traditione Maurolyci. Liber.}

\textbf{Propositio. 1.} Sub septentrionali polo habitantibus hemisphærium quidem mundi usquequaque idem apertum est unum alterum omnino idem occultum, Nec astrorum aliquod ipsis occidit, oriturve. Sed quæ in aperto sunt prorsus sunt
\note{la page s'arrête sur ces mots, laissant blanc le tiers inférieur ; la phrase continue en tête de la page suivante}

\note{page de droite de la vue 11, numérotée à l'encre « 10 » en haut à droite ; elle continue la page de gauche (« prorsus sunt » / « in conspectu »).}
in conspectu: at quæ in occulto usquequaque nunquam comparent.

Hæc propositio eadem ferme est cum quarta sphæræ Autolyci.

\textbf{Propositio. 2.} Sub æquinoctiali habitantibus omnia astra et oriuntur et occidunt, et æquali tempore super horizontem attollentur, horizontique subvehentur.
\note{sur cette page et la suivante, « Propositio » et le numéro sont en marge à droite}

Idem conclusit quinta propositio Autolyci.

\textbf{Propositio. 3.} Per omnem locum qui sub mediâ zonâ tropicis parallelis inclusâ signifer quotidie rectus erigitur horizonti.

Nam circulus æquatoris parallelus per verticem loci ductus binis in locis secat zodiacum. Quando igitur punctum sectionis alterutro counitur vertici, tunc zodiacus incedit per polos horizontis, et ideo per 20. primi Theodosii secat horizontem orthogonaliter. bis ergo fit hoc in uno ambitu. Habitantibus vero sub tropico semel in puncto solstitii in quo tangit zodiacus ipsum tropicum.

\textbf{Propositio 4.} Quorum vertex a polo tantum abit quantum tropicus distat ab æquinoctiali illis simul sex signa et oriuntur et occidunt.

Hoc est illis quorum vertex in artico vel antarctico circulo versatur. Nam cum poli zodiaci ferantur in periferiis talium circulorum Iam in uno ambitu semel polus counitur vertici hoc est polus zodiaci polo horizontis. Quare et zodiacus counitur horizonti, quæ counio fit in instanti et post illud instans statim zodiacus dispescitur bifariam ab horizonte Et ideo in instanti semicirculus eius oritur et reliquus semicirculus occidit.

\textbf{Propositio. 5.} Sub æquinoctiali habitantibus meridianus bifariam secabit super horizontem signiferi semicirculum, quando tactus tropicorum, et signiferi fuerint in horizonte, et tunc signifer rectus erit ad horizontem.

Tunc enim horizon incedens iam per polos tropici et puncta contactuum, per 6. 2. Sphæricorum Theodosij, ibit et per polos zodiaci, et ideo per 20. primi orthogonaliter eum secabit et per 21. primi zodiacus vicissim ibit per polos horizontis: per quos et meridianus Unde arcus tam meridiani quam zodiaci, ab horizontis polo ad horizontem recepti sunt quadrantes.

\textbf{Propositio. 6.} Sub æquinoctiali habitantibus signiferi semicirculi omnes in tempore æquali oriuntur similiter etiam

\page{12}
\note{double page ; page de gauche. La proposition 6, en bas de la page précédente (même feuillet), s'arrête sur « similiter etiam » et cette page commence avec la proposition 7 : la fin de l'énoncé et la démonstration de la proposition 6 ne sont pas écrites}

\textbf{Propositio. 7.} Habitantibus sub eodem parallelo stellæ neque simul oriuntur neque simul occidunt, sed quanto prius oriuntur orientaliori loco tanto prius occidunt.

Nam talium locorum horizontes propter æquas poli altitudines tangunt eosdem æquatoris parallelos quare per 18 secundi Sphæricorum Theodosij arcus ex parallelo quolibet semicirculis tam Orientalibus quam Occidentalibus interiecti similes sunt. Omnis igitur stella in loco Orientali per eumdem arcum anticipat ortum, et inde occasum, et idcirco per idem temporis intervallum.

\textbf{Propositio. 8.} Sub eodem meridiano habitantibus stellæ quæcumque sunt inter maximum semper apparentium parallelorum et æquinoctialem diutius super horizontem feruntur illis qui ad septemtrionem habitant quam illis qui ad meridiem. Et quanto prius oriuntur ad septemtrionem habitantibus, tanto posterius occidunt. Quæ vero astra inter maximum semper occultorum et æquinoctialem diutius super horizontem, apparent ad meridiem habitantibus, quam ad septentrionem, et quanto prius oriuntur his qui ad meridiem tanto posterius occidunt.

Nam eunti versus manifestum polum arcus diurnus astri eodem versus sub æquatore declinantis crescit et eunti versus occultum polum arcus diurnus eo declinantis etiam crescit Collatis autem arcubus utrimque hoc est ad ortum et occasum crescentibus constat reliqua pars propositi.

\textbf{Propositio. 9.} At si horizontes neque sub eodem parallelo neque sub eodem fuerint meridiano sequetur tantum arcuum super horizontem peractorum prædicto modo, inæqualitas. Non autem ortuum et occasuum anticipatio.

Constat sicut præmissa propter magis aut minus inclinatũ horizontem.\note{« inclinatũ » : ainsi écrit, avec un trait sur le « u ».}

\textbf{Propositio 10.} Sub utrolibet polo habitantibus sol semestri tempore super horizontem iugiter fertur, et tantumdem sub horizonte.

\note{page de droite de la vue 12, numérotée à l'encre « 11 » en haut à droite ; elle continue la page de gauche (proposition 10 / sa démonstration).}
Patet hoc per primam huius: quoniam Zodiaci semicirculus semper extat et semicirculus semper delitescit, qui semestri ferme spatio a sole peragitur. Neque hic motus differentia quam ingerit solis eccentricus consideranda venit. Supponitur enim solis motus æqualis. ubi de arcubus primi motus agitur.

\textbf{Propositio. 11.} A polo versus arcticum vel antarcticum procedentibus hæc iugis mora solis super horizontem aut sub horizonte minor fiet semestri tempore, minorque donec ad spatium viginti quatuor horarum sub arctico vel antarctico redigatur.

Nam horizon harum habitationum tangit duos æquatoris parallelos tropicis majores qui de zodiaco utrimque duas peripherias æquales abscindunt. Et peripheria quam parallelus semper apparens abscindit nunquam occidit eam vero quam semper occultus, nunquam oritur. In illa ergo quamdiu sol fuerit nunquam occidet. In hac nunquam orietur. Unde in illa tantus erit iugis dies. In hac tanta iugis nox: pro magnitudine scilicet peripheriæ tot scilicet mensium quot signa comprehenderit peripheria et tot dierum naturalium quot insuper gradus habuerit.

\textbf{Propositio 12.} Sub arctico vel antarctico habitantibus dies maxima fiet viginti quatuor horarum, et instans pro nocte, contra nox maxima 24 horarum et instans pro die. Cæteri arcus crescent et decrescent, usque ad æqualitatem æquinoctij.

Namque ibi habitantium horizon tangit duos tropicos, quos et zodiacus contingit. Quare sol in tropico super horizontem extantem constitutus, peragit pro die integrum ambitum. et pro nocte punctum ipsum contactus: Contra sol in tropico totaliter latenti positus circulum totum pro nocte describit et pro die solum contactus punctum, Inde decrescunt utrique arcus usque ad medium æqualitatis.

Et quoniam per 23 secundi Sphæricorum Theodosii coalterni arcus diurnus et nocturnus utrimque ab æquatore sumpti sunt æquales et perinde tam duo diurni quam duo nocturni coniugate recepti circulum integrant, propterea fit, ut in omni horizonte totum tempus diurnum in anno collectum simul sit semestre. hoc est anni dimidium et nocturnum similiter tantumdem. Sed sub polo continuum Cæteris vero habitationibus interpolatum.

FINIS.

\page{13}
\note{double page ; page de gauche. Un texte nouveau commence ; la page de droite de la vue 12 s'achève sur « FINIS. ».}
\section*{Euclidis Phænomena ex Traditione Maurolyci}

Quoniam stellæ circumferuntur æquali semper a nobis remotione propterea motum coeli circularem esse. Easque per parallelos æquidistantes deferri. Et quoniam quædam perpetuo feruntur supra terram quædam sub terrâ: quædam plus mora trahunt super terram: quædam plus sub terra et medio loco positæ æqualiter. Hinc mundum nonnisi figura sphærica esse, nec nisi circa axem et æqualiter volvi. Axis autem polorum unum extare, alterum sub terra delitescere. Horizon autem circulus vocetur qui definit spectatum hemisphærium. Meridianus qui per sphæræ horizontis polos incedat. Æquinoctialis vero maximus inter parallelos communem cum sphæra axem habentes. Zodiacum sive signiferum esse Solis orbitam a quo æquinoctialis oblique secatur. Tropici sunt duo paralleli æquales zodiacum tangentes. Ex quibus manifestum est horizontem circulum esse maximum quod maximum quemque semper secet per medium.
\note{une tache d'encre sur « et » après « sub terra »}

\textbf{Propositio 1.} Terra in medio mũdo est centrique fungitur officio.

Quoniam scilicet eadem dioptra spectamus duo signa opposita unum oriens alterum occidens. Et rursus alia duo signa opposita apud ortum et occasum: fit ut linea visualis in utraque observatione sit diametros zodiaci, et firmamenti, Terra igitur in sectione talium diametrorum cum sit in centro Zodiaci et perinde Mundi existet.

\textbf{Propositio 2.} In uno mundi ambitu qui per polos sphæræ circulus bis rectus erit ad horizontem.

Quia scilicet bis counitur in die cum meridiano, qui rectus est ad horizontem.

Zodiacus vero circulus ad meridianum bis erit rectus.

Quia scilicet polus Zodiaci in parallelo arctico delatus bis in die sistitur in meridiano.

Ad horizontem vero minime rectus erit, quoniam polus mundi fuerit extra tropicos.

Ibi enim Zodiacus nunquam transit per polos horizontis hoc est per verticem loci.

Si vero polus horizontis in tropico fuerit, Zodiacus circulus ad horizontem bis rectus erit.

Nam ibi parallelus æquatoris per polum horizontis incedens, binis in punctis secat zodiacum, quæ puncta singula semel in die sistentur in ipso polo horizontis in dicto parallelo delata, Bis igitur in die zodiacus horizontem orthogonaliter secabit per 20. 1. Sphæricorum elementorum Theodosij.

\textbf{Propositio 3.} Astrorum non errantium ortus occasusque efficientium unumquodque apud eadem horizontis puncta oritur, et occidit.

Nam parallelus in quo defertur astrum super axem, in uno semper situ circumducitur, et in iisdem semper punctis secat horizontem.

\note{page de droite de la vue 13, numérotée à l'encre « 12 » en haut à droite ; elle continue la page de gauche.}
\subsection*{Additio duarum propositionum.}
\note{le titre annonce deux propositions ajoutées ; la page en compte six, numérotées 4 à 9}

\textbf{Propositio. 4.} Astra in circulo per polos mundi ducto existentia simul oriuntur et simul occidunt in horizonte recto.

Nam talis circulus bis in die counitur horizonti recto.

\textbf{Propositio. 5.} Astra existentia in semicirculo orientali circuli tangentis maximum integre apparentium parallelorum quem tangit horizon obliquus, simul oriuntur in tali horizonte. Existentia vero in semicirculo reliquo simul occidunt in eodem.

Sic enim ille semicirculus semel in die counitur semicirculo orientali horizontis: ita hic occidentali. Unde denominantur.

\textbf{Propositio 6.} Astrorum in maximi circuli ambitu existentium maximumque integre apparentium non tangentis neque secantis, quæ prius oriuntur, prius occidunt: et quæ prius oriuntur prius occidunt.
\note{la seconde proposition répète la première telle qu'elle est écrite}

Nam ex talibus astris occidentalibus prius oritur et prius occidit. Ducto enim semicirculo orientali tangente maximum parallelorum integre apparentium per astrum occidentalius relinquitur astrum reliquum ad orientem. Similiter ducto semicirculo occidentali. Constat ergo propositum: cum tales semicirculi representent semicirculos horizontis.
\note{« occidentalibus » : ainsi écrit, là où la construction attend un singulier ; plus bas, « per astrum occidentalius ».}

\textbf{Propositio. 7.} Astrorum in maximi orbis ambitu, qui maximum integre apparentium secat existentium, quæ septentrioni propius prius oriuntur, posterius vero occidunt.

Ductis enim per astrum a septentrione remotius semicirculis majoribus maximum integre apparentium utrimque tangentibus: relinquetur astrum reliquum in medio periph\add{e}riarum. Unde palam fit ipsum astrum reliquum prius oriri, et posterius occidere. Sed Euclides loquitur respectu situs nostri. Nam apud nostros antoecos idem dicendum est de astro, quod illi polo propinquius est.

\textbf{Propositio. 8.} In zodiaco sive æquinoctiali, sive quovis alio maiori circulo astra ex diametro posita coniugate oriuntur et occidunt.

Nam quovis diameter cuiuslibet majoris circuli est et mundi diameter: cuius extremorum altero exoriente alterum occidit, et e contrario.

\textbf{Propositio. 9.} Zodiacus circulus per omnem horizontis locum inter circulos tropicos oritur, quando maximus integre
\note{la page s'arrête ici, au milieu de l'énoncé}

\page{14}
\note{double page ; page de gauche. Elle continue la page de droite de la vue 13 (« quando maximus integre » / « apparentium »).}
apparentium non \uncertain{minor} fuerit circulo tropico.
\note{« minor » est surchargé d'encre}

Hoc est in illo horizonte, cuius vertex est in circulo artico, vel inter ipsum et polum, ortus zodiaci fit per totum semicirculum horizontis \uncertain{orientalem}: occasus per totum semicirculum horizontis occidentalem: quandoquidem totus horizon iacet intra tropicos.
\note{« orientalem » : le milieu du mot est couvert d'une tache d'encre ; seuls « or » et « entalem » se lisent}

\textbf{Propositio 10.} Signa non apud æqualia horizontis segmenta oriuntur et occidunt, in maximis enim quæ ad æquinoctialem, in minoribus autem quæ hæc sequuntur in minimis quæ ad tropicos apud æqualia porro quæ ab æquinoctiali circulo æqualiter distant.

Ductis enim per limites signorum Zodiaci parallelis hinc inde ab æquinoctiali: periferiæ horizontis interceptæ tam ad ortum quam ad occasum, ab æquinoctiali versus tropicos ordinata successione decrescunt. Ut infert propositio. Omnis enim arcus zodiaci apud peripherias horizontis suis parallelis interceptas oritur et occidit. Hoc autem ostendit Theodosius in propositionibus 3. 5. 7 et 9. tertii Sphæricorum Menelai aut in 46 2. Quod intelligitur tam in horizonte recto, quam in obliquo, quamvis in obliquo peripheriæ dictæ horizontis sint majores.
\note{« Theodosius … Spharicorum Menelai » est écrit ainsi ; « 46 2 » semble renvoyer à la proposition 46 d'un livre II}

\textbf{Propositio 11.} Zodiaci semicirculos non ab eodem parallelo exorsos neque æquali tempore exoriri, sed in plurimo, qui cum cancro, eoque minori, qui inde remotius exordium sumpserint, in minimo tandem qui cum Capricorno. Quicumque autem in eodem parallelo initium habuerint: æquis temporum intervallis exortum facere.

Constat hoc propositio apertissime si conferuntur arcus diurni semicirculorum Zodiaci principiis debiti. Cum talibus enim arcubus oriuntur ipsi semicirculi. et pro occasu semicirculorum conferuntur arcus nocturni, qui semicirculorum initijs respondent: quamvis autor de occasu non faciat mentionem sed ultra æquinoctialem pro signo in propositione expresso sume signo opposita. Autor enim respexit ad situm nostrum.

\textbf{Propositio 12.} Si in zodiaco bini semicirculi communem quandam circumferentiam habentes in diverso tempore oriantur relicti arcus diverso etiam tempore orientur et in eodem excessu, Si autem in zodiaco bini

\note{page de droite de la vue 14, numérotée à l'encre « 13 » en haut à droite ; elle continue la page de gauche (« Si autem in zodiaco bini » / « communem arcum »).}
communem arcum habentes æquis temporibus oriantur relictæ quoque peripheriæ temporibus æquis orientur.

Nam subtracto arcu communi, subtrahitur etiam \uncertain{eius} tempus et ideo relicta tempora erunt aut in eodem excessu inæqualia. In quo scilicet tempora semicirculorum sunt aut æqualia, si tempora semicirculorum fuerint æqualia.
\note{« eius » rend un mot abrégé, surmonté d'un trait, qui se lit à peu près « eoc »}

\textbf{Propositio. 13.} Zodiaci æqualium et ex opposito circumferentiarum in quo tempore altera oritur, reliqua occidit et in quo altera occidit, reliqua \struck{occidit} oritur.

Nam per 8 huius talium arcuum limites exeuntes ex diametro coniugate oriuntur, et occidunt, hoc est \struck{in} uno oriente alter occidit et e contrario. Et ideo quo tempore oritur interceptorum arcuum unus reliquus cooccidit: et e contrario.

\subsection*{Additio trium propositionum.}

\textbf{Propositio. 14.} Similium horizontes semicirculi similes parallelorum periferias includunt: et ideo quodlibet astrum ad horizontem ex his orientalem per unum temporis intervallum anticipat tam ortum quam occasum ac coeli mediationem.

Similes horizontes sunt qui aut recti sunt aut eiusdem latitudinis: Qui autem eiusdem latitudinis sunt tangunt eosdem parallelos, quorum alter maximus integre apparentium est, alter maximus integre occultorum. Hæc ergo propositio quo ad rectos horizontes ostenditur in Sphæricorum Theod. 14. secundi: quo ad autem obliquos in 8 eiusdem. Ut si inter duos horizontum sive rectorum sive unius latitudinis obliquorum semicirculos orientales intersit arcus æquatoris 30 graduum: Iam inter eosdem ex quolibet parallelo totidem gradus intercipientur. Et perinde omne astrum magis orientale per duas horas prævertet tam ortum quam occasum, quam et coeli mediationem. Quare constat aperte corollarium.

\textbf{Propositio 15.} Similium horizontum semicirculi orientales una cum periferijs intercipiunt æquatoris arcus coorientes, occidentales aut cooccidentales ad quemlibet talium horizontum.

Manente enim fixo horizontum talium uno sphæra revoluta cæterorum horizontum semicirculi cooriuntur ei, et perinde zodiaci peripheriæ ante motum interceptæ cooriuntur aut coocciduunt cum arcubus æquatoris simul interceptis.

\textbf{Propositio 16.} Peripheriæ zodiaci æquales ad rectum horizontem non æquis temporibus oriuntur neque occidunt sed in maximo quæ sunt ad tropicorum contactus: In minori autem, quæ has subsequuntur, In minimis vero quæ ad æquinoctialem
\note{la page s'arrête ici ; la phrase continue sur la page suivante}

\page{15}
\note{double page ; page de gauche. Elle continue la page de droite de la vue 14 (« quæ ad æquinoctialem » / « æqualibus porro »).}
æqualibus porro quæ ab æquinoctii puncto æqualiter distant.

Exempli gratia sumantur in Zodiaco tria signa Aries Taurus et Gemini. Aio quod ex his in sphæra recta Gemini in maximo: Taurus in minori: Aries in minimo tam oritur quam occidit tempore Ducantur enim a polo mundi tres semicirculi horizontum rectorum per limites talium signorum: Iam tales semicirculi abscindent de æquinoctiali arcus inæquales: quorum maximus erit, qui remotissimus a sectione Zodiaci et æquinoctialis scilicet qui cum Geminis intercipitur: minor qui cum Tauro, minimus qui cum ariete: per 4. et 8. 3\textsuperscript{i} Sphæricorum Theodo. et 46 secundi Sphæricorum Menelai. Sed per præcedentem tales æquatoris arcus cooriuntur sive coocciduunt, cum signis ipsis interceptis Igitur ex his Gemini in maximo Taurus in minori. Aries in minimo oritur et occidit tempore. Quod autem æque ab æquinoctio remota signa æquis temporibus oriuntur et occidunt. constat quoniam cum æquis arcubus æquatoris oriuntur et occidunt, \struck{constat quoniam} et id propter æquilatera invicem sphæralia triangula per 25. primi Sphæricorum Menelai.

\subsection*{Corollarium.}

Hinc manifestum est quod in sphærâ recta quatuor signa Gemini Cancer Sagittarius et Capricornus, in maximis autem quatuor autem Taurus Leo Scorpius et Aquarius in minoribus quatuor demum Pisces, Aries, Virgo et libra in minimis et invicem æqualibus oriuntur et occidunt temporibus

\textbf{Propositio 17.} Semicirculi Zodiaci qui cum Cancro æquales circumferentiæ non æquis temporibus occidunt Sed in maxima quæ sunt ad Tropicorum contactus: In minori autem quæ has subsequuntur, In minimis vero quæ ad æquinoctialem, æqualibus porro quæ ab æquinoctij puncto æqualiter distant et occidunt.

Per limites trium signorum Cancri Leonis et Virginis ducantur tres semicirculi horizontum obliquorum eiusdem latitudinis occidentales, et proinde tangentes eumdem æquatoris parallelos: Nam tales semicirculi abscindent ex æquatore arcus inæquales, quorum maximus erit qui remotissimus a sectione æquinoctij, scilicet qui cum Cancro intercipitur. Minor qui cum Leone, minimus qui cum Virgine per 6. et 10 tertij Sphæricorum Theod. et 46. Sphæricorum Menelai. Sed per antepramissam talia signa coocciduunt cum arcubus æquatoris interceptis. Igitur ex his Cancer in maximo, Leo in minori. Virgo in minimo occidunt tempore.
\note{« et occidunt » à la fin de l'énoncé de la proposition 17 est écrit ainsi}

\note{page de droite de la vue 15, numérotée à l'encre « 14 » en haut à droite ; elle continue la page de gauche.}
Quod autem signa æqualiter ab æquinoctio remota æquis temporibus oriuntur et occidunt, constat per 25. primi Sphæricorum Menelai, propter æquilatera invicem sphæralia triangula. Verum in horizontib\textsuperscript{9} ultra æquatorem pro Cancro Leone et Virgine substitue Capricornum Aquarium, et Pisces.

\subsection*{Corollarium.}

Hinc patet quod in horizonte nostro obliquo, dua signa Cancer et Sagittarius in maximis: Leo et Scorpius in minoribus: Virgo et Libra in minimis et invicem æqualibus occidunt temporibus.

\textbf{Propositio. 18.} Semicirculi Zodiaci qui cum Capricorno æquales periferiæ nequaquam æquis oriuntur temporibus In maximo quidem, quæ ad Tropicorum contactus In minori autem quæ has subsequuntur: In minimis vero quæ ad æquinoctialem. æqualibus porro, quæ ab æquinoctii puncto æqualiter distant oriuntur et occidunt.

Ostensum est in præcedenti Cancrum in maximo, in minori Leonem, in minimo Virginem occidere. Igitur per 13. præcedentem his opposita signa scilicet Capricornus in maximo, Aquarius in minori: Pisces in minimo orietur. quod est propositum. Unde et æqualitas ortuum in signis æque ab æquinoctio remotis sequetur. Sed in regionibus ultra æquatorem: quoniam mutatur polus manifestus. Commutanda sunt signa.

\subsection*{Corollarium.}

Constabit igitur hic similiter quod in his obliquis horizontibus Capricornus et Gemini in maximis, Aquarius et Taurus in minoribus, Pisces et Aries in minimis et invicem æqualibus oriuntur temporibus.
\note{« Gemini » et « Taurus » sont écrits là où le parallèle avec le corollaire précédent attendrait « Sagittarius » et « Scorpius »}

Sic Euclidis Phoenomena et Menelai et Sphæricorum Theodosii adiumentis ubi multo facilius quam autor fecerat demonstravimus, unde patet quanta sit non solum commoditas sed etiam necessitas Sphæricorum in astronomicis rudimentis demonstrandis.

\subsection*{Habitationum collatio.}

Perioeci dicuntur sub eodem parallelo habitantes quasi circumcola. His eadem est et eiusdem poli celsitudo, æquales simul arcus tam diurni quam nocturni, eadem astrorum et signorum apparitio, occultatio ortus et occasus. Simul habent æstatem. simul hiemem. simul ver simul autumnum. Easdem meridianas umbras. Per idem intervallum anticipat tam ortum quam occasum.

Antoeci sunt qui sub æqualibus et oppositis parallelis habitant quasi contracola. His æqualis est sed diversorum polorum celsitudo. Æquales sed oppositorum punctorum arcus tam diurni quam nocturni. Eadem sed oppositorum

\page{16}
\note{double page ; page de gauche. Elle continue la page de droite de la vue 15 (« Eadem sed oppositorum » / « astrorum »).}
astrorum et signorum apparitio, occultatio, ortus et occasus Quando his fit æstas, illis hyems Quando his fit ver, illis autumnus. Æquales habent sed in oppositis illis locis et in diversum proiectas meridianas umbras.

Antichthones sive Antipodes sunt qui non solum sub æqualibus et oppositis parallelis habitant sed etiam in locis per diametrum oppositis atque contrarijs invicem pedibus terram calcantes. Itaque omnia quæ antoecis accidunt etiam ad antipodes omnino spectant. Sed Antipodum hoc est proprium quod habent communem horizontem, et hemisphæria diversa ac vertices diversos, et quicquid his oritur illis occidit. Quicquid his extat illis delitescit. Quicquid his ascendit illis descendit.

Amphiscij sunt qui intra tropicos habitant his enim umbra meridiana quandoque dextrorsum quandoque sinistrorsum soli orienti proijcitur. Qui etiam duas æstates et totidem hiemes habent.

Periscij sunt qui sub polo. his enim umbra per totum horizontis ambitum sole similiter circumlato defertur.

Antiscij sunt perioeci quibus quidem æqualis et eodem versus meridiana umbra flectitur.

Heteroscij sunt qui et Antoeci quibus in diversum proijcitur umbra meridiana, Possunt intelligi heteroscij quibus altera umbrarum tantum hoc est versus unum polorum in meridie proijcitur. Ut pote extra Tropicos habitantibus.

\subsection*{DE Astrorum fulsionibus.}

Intervallum quod stellam in horizonte positam a sole in prima vel postrema fulsione disterminat, neque in Zodiaco est computandum, neque idem est in omnibus astris sicut autolycus in Phænomenis supposuit, ac fortasse putavit. Sed cum talis distantia respectu horizontis consideranda sit, Iam in circulo per polos horizontis ducto capienda est: Et in stellis luminosioribus, quæ prius apparent et posterius evanescunt minor esse debet. Unde Venus planetarum lucidissima minimum postulat intervallum. Lege super his 8 magnæ constructionis Ptolemæicæ. Putant etiam nonnulli Lunam minus 14 gradibus a sole remotam apparere non posse. Quod falsum est. Nam intra multo minus spatium sæpe visa est, neque spatium in zodiaco mensurandum. Item non recte Plinius ex intervallo tali arguit lunam cæteris astris minorem. Nam intervallum maius arguit astrum minus luminosum: hoc est minus in apparentia non in vera magnitudine. Intervallum quidem quod solem ab horizonte removet in prima postremave fulsione veneris est quinque graduum, Ptolemæi testimonio. Hinc servata luminis proportione Ut lunæ primum postremove apparenti satis est trium graduum interstitium. Nam cum tantum interstitium sit sexagesima

\note{page de droite de la vue 16, numérotée à l'encre « 15 » en haut à droite ; elle continue la page de gauche (« sexagesima » / « pars semicirculi »).}
pars semicirculi, ut tunc luna adepta sit de plenilunio toti\note{« toti » est traversé d'un trait oblique, peut-être une biffure.} lumini sexagesimam. Sed plena Luna centupla ferme est venereæ faciei, \struck{cuius} illius diameter ad huius diametrum decupla ponatur ab Albategnio, Et ideo illud lunæ corniculatum lumen tenue ad lumen veneris erit sicut quinque ad 3. Igitur si Venus ad visionem postulat quinque graduum interstitium: Lunæ satis erit trium graduum intercapedo. quandoquidem intervalla luminibus reciproca esse debent. Unde si coniunctio luminarium fiat iuxta meridiem in diebus æstivis et expeditis horizontibus eodem die luna videri poterit matutina et vespertina.

\subsection*{Tabella sinus recti.}
\note{le titre est écrit à droite, au bout d'un trait qui part de la dernière ligne du texte. La table occupe le reste de la page : deux blocs de quatre colonnes, degrés 1 à 45 et 46 à 90, séparés de cinq en cinq par des filets ; les intitulés des colonnes sont écrits au pied, « gradus », « sinus », « Dria » (abrégé, avec un tilde, pour « Differentia »), « 6\textsuperscript{ma} » à gauche et « gradus », « sinus », « Differentia », « 60\textsuperscript{ma} » à droite. On la donne ici en trois tranches de quinze lignes, sans rien changer à l'ordre des lignes. Aux lignes 46 et 47, la dernière colonne porte un point au lieu des deux-points. Plusieurs chiffres sont repassés ou surchargés (10 : 28:42 ; 24 : 40674 ; 35 : 23:59 ; 39 : 22:46, le 6 ajouté au-dessus ; 48 : 74315 ; 59 : 85717 ; 63 : 89101) ; on donne la lecture finale.}

\[
\begin{array}{r|r|r|r||r|r|r|r}
1 & 1745 & 1745 & 29{:}5 & 46 & 71934 & 1223 & 20.23 \\
2 & 3490 & 1745 & 29{:}5 & 47 & 73135 & 1201 & 20.1 \\
3 & 5234 & 1744 & 29{:}4 & 48 & 74315 & 1180 & 19{:}40 \\
4 & 6976 & 1742 & 29{:}2 & 49 & 75471 & 1156 & 19{:}16 \\
5 & 8716 & 1740 & 29{:}0 & 50 & 76604 & 1133 & 18{:}53 \\
\hline
6 & 10453 & 1737 & 28{:}57 & 51 & 77715 & 1111 & 18{:}31 \\
7 & 12187 & 1734 & 28{:}54 & 52 & 78801 & 1086 & 18{:}6 \\
8 & 13917 & 1730 & 28{:}50 & 53 & 79864 & 1063 & 17{:}43 \\
9 & 15643 & 1726 & 28{:}46 & 54 & 80902 & 1038 & 17{:}18 \\
10 & 17365 & 1722 & 28{:}42 & 55 & 81915 & 1013 & 16{:}53 \\
\hline
11 & 19081 & 1716 & 28{:}36 & 56 & 82904 & 989 & 16{:}29 \\
12 & 20791 & 1710 & 28{:}30 & 57 & 83867 & 963 & 16{:}3 \\
13 & 22495 & 1704 & 28{:}24 & 58 & 84805 & 938 & 15{:}38 \\
14 & 24192 & 1697 & 28{:}17 & 59 & 85717 & 912 & 15{:}21 \\
15 & 25882 & 1690 & 28{:}10 & 60 & 86603 & 886 & 14{:}46 \\
\end{array}
\]

\[
\begin{array}{r|r|r|r||r|r|r|r}
16 & 27564 & 1682 & 28{:}2 & 61 & 87462 & 859 & 14{:}19 \\
17 & 29237 & 1673 & 27{:}53 & 62 & 88295 & 833 & 13{:}53 \\
18 & 30902 & 1665 & 27{:}45 & 63 & 89101 & 806 & 13{:}26 \\
19 & 32557 & 1655 & 27{:}35 & 64 & 89879 & 778 & 12{:}58 \\
20 & 34202 & 1645 & 27{:}25 & 65 & 90631 & 752 & 12{:}32 \\
\hline
21 & 35837 & 1635 & 27{:}15 & 66 & 91355 & 724 & 12{:}4 \\
22 & 37461 & 1624 & 27{:}4 & 67 & 92051 & 696 & 11{:}36 \\
23 & 39073 & 1612 & 26{:}52 & 68 & 92718 & 667 & 11{:}7 \\
24 & 40674 & 1601 & 26{:}41 & 69 & 93358 & 640 & 10{:}40 \\
25 & 42262 & 1588 & 26{:}28 & 70 & 93969 & 611 & 10{:}11 \\
\hline
26 & 43837 & 1575 & 26{:}15 & 71 & 94552 & 583 & 9{:}43 \\
27 & 45399 & 1562 & 26{:}2 & 72 & 95106 & 554 & 9{:}14 \\
28 & 46947 & 1548 & 25{:}48 & 73 & 95631 & 525 & 8{:}45 \\
29 & 48481 & 1534 & 25{:}34 & 74 & 96126 & 495 & 8{:}15 \\
30 & 50000 & 1519 & 25{:}19 & 75 & 96593 & 467 & 7{:}47 \\
\end{array}
\]

\[
\begin{array}{r|r|r|r||r|r|r|r}
31 & 51504 & 1504 & 25{:}4 & 76 & 97030 & 437 & 7{:}17 \\
32 & 52992 & 1488 & 24{:}48 & 77 & 97437 & 407 & 6{:}47 \\
33 & 54464 & 1472 & 24{:}32 & 78 & 97815 & 378 & 6{:}18 \\
34 & 55919 & 1455 & 24{:}15 & 79 & 98163 & 348 & 5{:}48 \\
35 & 57358 & 1439 & 23{:}59 & 80 & 98481 & 318 & 5{:}18 \\
\hline
36 & 58779 & 1421 & 23{:}41 & 81 & 98769 & 288 & 4{:}48 \\
37 & 60182 & 1403 & 23{:}23 & 82 & 99027 & 258 & 4{:}18 \\
38 & 61566 & 1384 & 23{:}4 & 83 & 99255 & 228 & 3{:}48 \\
39 & 62932 & 1366 & 22{:}46 & 84 & 99452 & 197 & 3{:}17 \\
40 & 64279 & 1347 & 22{:}27 & 85 & 99620 & 168 & 2{:}48 \\
\hline
41 & 65606 & 1327 & 22{:}7 & 86 & 99756 & 139 & 2{:}16 \\
42 & 66913 & 1307 & 21{:}47 & 87 & 99863 & 107 & 1{:}47 \\
43 & 68200 & 1287 & 21{:}27 & 88 & 99939 & 76 & 1{:}6 \\
44 & 69466 & 1266 & 21{:}6 & 89 & 99985 & 46 & 0{:}46 \\
45 & 70711 & 1245 & 20{:}45 & 90 & 100.000 & 15 & 0{:}15 \\
\hline
\text{gradus} & \text{sinus} & \text{Dria} & 6^{\text{ma}} & \text{gradus} & \text{sinus} & \text{Differentia} & 60^{\text{ma}}
\end{array}
\]

\note{ligne 4, deuxième colonne : le troisième chiffre est formé presque comme un 1 (« 1712 ») ; il est lu 4, ce que 29:2 dans la colonne suivante appuie ; lecture douteuse. Trois valeurs de la dernière colonne ne s'accordent pas avec la différence qu'elles divisent par 60 et restent telles qu'écrites : ligne 59, 912 et 15:21 (912/60 donne 15:12) ; ligne 86, 139 et 2:16 (99756 − 99620 = 136, et 136/60 donne 2:16) ; ligne 88, 76 et 1:6 (76/60 donne 1:16).}

\page{17}
\note{double page ; page de gauche, une table : la « Tabula foecunda », qui fait suite à la table de sinus de la vue 16.}
\section*{Tabula foecunda}
\note{sous ce titre, une table en deux blocs, degrés 1 à 45 et 46 à 90, colonnes intitulées en tête « gr. », « umbra », « Dria » (abrégé pour « Differentia »), « 6\textsuperscript{ma} » ; au second bloc, « g. » pour « gr. ». Filets de cinq en cinq lignes. Donnée ici en trois tranches de quinze lignes. Le point ou les deux-points de la dernière colonne sont gardés tels qu'écrits.}

\[
\begin{array}{r|r|r|r||r|r|r|r}
\text{gr.} & \text{umbra} & \text{Dria} & 6^{\text{ma}} & \text{g.} & \text{umbra} & \text{Dria} & 6^{\text{ma}} \\
\hline
1 & 1745 & 1745 & 29.5 & 46 & 103551 & 3551 & 59{:}11 \\
2 & 3492 & 1747 & 29.7 & 47 & 107236 & 3685 & 61.25 \\
3 & 5241 & 1749 & 29.9 & 48 & 111062 & 3826 & 63{:}46 \\
4 & 6992 & 1751 & 29{:}11 & 49 & 115037 & 3975 & 66.15 \\
5 & 8748 & 1756 & 29{:}16 & 50 & 119197 & 4160 & 69.20 \\
\hline
6 & 10510 & 1762 & 29{:}22 & 51 & 123491 & 4294 & 71{:}34 \\
7 & 12278 & 1768 & 29{:}28 & 52 & 127994 & 4503 & 75{:}3 \\
8 & 14053 & 1775 & 29.35 & 53 & 132704 & 4710 & 78{:}30 \\
9 & 15838 & 1785 & 29{:}45 & 54 & 137639 & 4935 & 82{:}15 \\
10 & 17633 & 1795 & 29{:}55 & 55 & 142813 & 5174 & 86{:}14 \\
\hline
11 & 19439 & 1806 & 30{:}6 & 56 & 148253 & 5440 & 90{:}40 \\
12 & 21256 & 1817 & 30.17 & 57 & 153987 & 5734 & 95{:}34 \\
13 & 23087 & 1831 & 30{:}31 & 58 & 160033 & 6046 & 100.46 \\
14 & 24932 & 1845 & 30{:}45 & 59 & 166429 & 6396 & 106{:}36 \\
15 & 26794 & 1862 & 31{:}2 & 60 & 173205 & 6776 & 112{:}56 \\
\end{array}
\]

\[
\begin{array}{r|r|r|r||r|r|r|r}
16 & 28674 & 1880 & 31{:}20 & 61 & 180405 & 7200 & 120{:}0 \\
17 & 30573 & 1899 & 31{:}39 & 62 & 188073 & 7668 & 127{:}48 \\
18 & 32492 & 1919 & 31{:}59 & 63 & 196261 & 8188 & 136{:}28 \\
19 & 34433 & 1941 & 32{:}21 & 64 & 205030 & 8769 & 146{:}9 \\
20 & 36396 & 1963 & 32{:}43 & 65 & 214451 & 9421 & 157{:}1 \\
\hline
21 & 38387 & 1991 & 33{:}11 & 66 & 224603 & 10152 & 169{:}12 \\
22 & 40402 & 2015 & 33{:}35 & 67 & 235585 & 10982 & 183{:}2 \\
23 & 42448 & 2046 & 34{:}6 & 68 & 247509 & 11924 & 196{:}44 \\
24 & 44522 & 2074 & 34{:}34 & 69 & 260509 & 13000 & 216{:}4 \\
25 & 46631 & 2109 & 35{:}9 & 70 & 274747 & 14238 & 237{:}18 \\
\hline
26 & 48772 & 2141 & 35{:}41 & 71 & 290421 & 15674 & 261{:}14 \\
27 & 50952 & 2180 & 36{:}20 & 72 & 307768 & 17347 & 289{:}7 \\
28 & 53170 & 2218 & 36{:}58 & 73 & 327084 & 19316 & 321{:}56 \\
29 & 55432 & 2262 & 37{:}42 & 74 & 348742 & 21658 & 360{:}58 \\
30 & 57735 & 2303 & 38{:}23 & 75 & 373205 & 24463 & 407{:}43 \\
\end{array}
\]

\[
\begin{array}{r|r|r|r||r|r|r|r}
31 & 60086 & 2351 & 39{:}11 & 76 & 401078 & 27873 & 464{:}33 \\
32 & 62486 & 2400 & 40.0 & 77 & 433148 & 32070 & 534{:}30 \\
33 & 64940 & 2454 & 40{:}54 & 78 & 470453 & 37305 & 621{:}45 \\
34 & 67452 & 2512 & 41{:}51 & 79 & 514455 & 44002 & 733{:}22 \\
35 & 70022 & 2570 & 42{:}50 & 80 & 567128 & 52673 & 877{:}53 \\
\hline
36 & 72654 & 2632 & 43{:}52 & 81 & 631375 & 64247 & 1070{:}47 \\
37 & 75356 & 2702 & 45{:}2 & 82 & 711537 & 80162 & 1336{:}2 \\
38 & 78129 & 2773 & 46{:}13 & 83 & 814435 & 102898 & 1714{:}18 \\
39 & 80978 & 2849 & 47{:}29 & 84 & 951436 & 137061 & 2283{:}21 \\
46 & 83909 & 2931 & 48{:}51 & 85 & 1143005 & 191569 & 3192{:}49 \\
\hline
41 & 86929 & 3020 & 50{:}20 & 86 & 1430067 & 287062 & 4784{:}22 \\
42 & 90040 & 3111 & 51{:}51 & 87 & 1908113 & 478046 & 7967{:}26 \\
43 & 63252 & 3212 & 53{:}32 & 88 & 2863625 & 955512 & 15925{:}12 \\
44 & 96571 & 3319 & 55{:}19 & 89 & 5728993 & 2865370 & 47756{:}10 \\
45. & 100000 & 3429 & 57{:}9 & 90. & \text{Infinitum} & \text{Infinitum} & \text{Infinitum} \\
\end{array}
\]

\note{tel qu'écrit : ligne 40, le numéro du degré est « 46 » ; ligne 43, l'ombre est « 63252 », entre 90040 et 96571 ; ligne 34, 2512 et 41:51 ; ligne 69, « 216: 4 », sans chiffre après le 4. Ligne 4, dernière colonne, « 11 » est une lecture douteuse (le second chiffre ressemble à un b). Chiffres surchargés ou repassés : 16 (28674), 38 (46:13), 68 (numéro du degré), 74 (21658), 75 (24463).}

\[
\begin{array}{lcl}
89.\;15 & \ldots\ldots & 7638998 \\
89{:}30 & \ldots\ldots & 11458872 \\
89\;\;45 & \ldots\ldots & 22918163 \\
89{:}\,55 & \ldots\ldots & 68754439 \\
89{:}\,59 & \ldots\ldots & 343772546.
\end{array}
\]
\note{ces cinq lignes sont écrites sous la table, hors des colonnes, à la même encre ; la page est blanche au-dessous}

\note{page de droite de la vue 17, numérotée à l'encre « 16 » en haut à droite. Elle ne continue pas la page de gauche, une table : elle s'ouvre au milieu d'une règle pour trouver l'ascension droite d'une étoile, dont le début est au bas de la page de gauche de la vue 19 (« Nam in primis » / « considerandum est »). Voir la note de la vue 19.}

considerandum est in quo ex quatuor quadrantibus Zodiaci sit stellæ locus, an scilicet in primo quadrante a principio arietis ad finem Geminorum: an in 2\textsuperscript{o} a principio Cancri ad finem Virginis: an in 3\textsuperscript{o} a principio libræ ad finem Sagittarii, an denique in quarto a principio Capricorni ad finem piscium. Item cuius sit nominis stellæ declinatio scilicet septentrionalis an meridionalis Si enim stellæ locus fuerit in primo vel 2\textsuperscript{o} quadrante cum declinatione septentrionali aut in tertio vel quarto quadrante cum declinatione australi: tunc differentia transitus stellæ per coeli medium subtrahenda est ab arcu tertio, vel arcus a differentia minor scilicet de majori: Sic enim supererit arcus æquinoctialis inter punctum æquinoctii proximum et circulum declinationis stellæ clausus Quare si locus stellæ in primo quadrante fuit et differentia minor arcu tertio, tunc talis arcus æquinoctialis erit ascensio quæsita Si autem arcus tertius minor quam differentia, tunc dictus arcus æquinoctialis de toto circulo surreptus relinquet ascensionem quæsitam. \uncertain{At} vero locus stellæ in secundo quadrante fuit, et differentia minor arcu tertio tunc arcus iste demptus a semicirculo residuabit ascensionem. Si arcus \struck{arcu} tertius minor quam differentia tunc arcus iunctus cum semicirculo conflabit quæsitam ascensionem. Quod si locus stellæ in tertio quadrante extiterit, et differentia minor arcu tertio, arcus relictus iunctus ad hunc cum semicirculo conficiet ascensionem. Si arcus tertius minor quam differentia arcus relictus subtractus a semicirculo \struck{conficiet} relinquet ascensionem. \struck{Si arcus tertius minor quam differentia, tunc arcus relictus de toto arcus relictus iunctus}
\note{« relinquet » est écrit dans l'interligne au-dessus de « conficiet » biffé. Dans la ligne biffée qui suit, le trait s'arrête avant le dernier mot, « iunctus », qui reste peut-être hors de la rature.}
Si tandem locus stellæ in quarto quadrante fuerit, et dria\note{abréviation de « differentia », écrite « dria » avec un trait.} minor arcu tertio, tunc arcus relictus de toto circulo \uncertain{relinquet} ascensionem, Si arcus tertius fuerit minor quam differentia tunc arcus relictus pro ascensione tenendus est.
\note{« relinquet » : les premières lettres sont surchargées d'une tache}

At si locus stellæ fuerit in quadrante Zodiaci primo vel secundo cum declinaẽ\note{abréviation de « declinatione ».} australi aut in quadrante tertio vel quarto cum declinatione septentrionali, tunc differentia transitus stellæ per medium coeli iungenda est cum arcu tertio ut conflatur arcus æquinoctialis inter punctum æquinoctii propius et circulum declinationis stellæ comprehensus. Verum tunc si locus stellæ sit in primo quadrante is arcus est ipsa recta ascensio. Si autem locus sit in secundo quadrante tunc arcus demptus de semicirculo relinquet ascensionem quæsitam. Si locus sit in 3\textsuperscript{o} quadrante tunc arcus idem iunctus cum semicirculo conflabit ascensionem. Si demum locus in quarto quadrante fuit arcus de toto circulo sublatus residuabit quæsitam stellæ ascensionem. Denique attendendum quod quoniam locus stellæ est in principio arietis \struck{tun} aut libræ, tunc arcus primus et tertius nulli sunt quare agendum est cum solo arcu secundo, qui tunc habet gradus 66. m. 30. complementum scilicet maximæ
\note{« 66 » : le premier chiffre est le 6 en forme de « b » de cette main ; 66 gr. 30 m. est aussi le second arc de la ligne 1 de la table de la vue 20. La page s'arrête sur « maximæ », sans « declinationis » ; la suite est en tête de la page de gauche de la vue 18.}

\page{18}
\note{double page ; page de gauche. Elle continue la page de droite de la vue 17 (« complementum scilicet maximæ » / « declinationis »).}
declinationis \struck{Zodiaci et latitudinis stellæ si fuerint eiusdem nominis cum ambobus} et pro argumento declinationis tenenda est stellæ latitudo, Et tunc si locus stellæ fuerit in principio Libræ cum Arietis cum declinatione australi arcus inventus ut iacet erit ipsa ascensio. Si cum declinatione septentrionali ablatus de toto \struck{semi}circulo residuabit ascensionem quæsitam. Si autem stellæ locus sit in principio libræ cum declinatione septentrionali, arcus sæpe memoratus cum semicirculo eamdem residuabit. Porro si stellæ cuiuspiam alterum solstitialium punctorum locus in Zodiaco contingat, tunc declinatio stellæ habetur per adgregationem maximæ declinationis Zodiaci et latitudinis stellæ, si fuerint eiusdem nominis, cum ambobus vel cum majori. Circulus enim latitudinis et declinationis ibi unus et idem est, ipse scilicet colurus solstitialis. Item stellæ in solstitio æstivo constitutæ ascensio recta semper est circuli quadrans. In reliquo In reliquo positæ dodrans.
\note{les deux premières lignes sont biffées de « Zodiaci » à « ambobus ». « Libræ cum Arietis » est écrit ainsi. « In reliquo » est écrit deux fois de suite.}
Exempli gratia stellæ aldebaram declinatio fuit gr. 15. m. 43. \struck{Et complementum arcus secundi scilicet graduum} arcus secundus graduum 79. m. 31. arcus tertius gr. 64. m. 47. Igitur cum declinatione gr. 15. m. 43 et compl\textsuperscript{to} arcus secundi scilicet g. 10. m. 27. ex tabula foecunda elicio duos numeros scilicet 28140 atque 18506 quibus invicem multiplicatis produco 520758840. \uncertain{Demo} abicio quinque dextimas figuras et supersunt 5208. cuius sinus arcus est paulo minor quam graduum 3. differentia scilicet transitus stellæ per coeli medium. Vel sic cum \uncertain{dicta} declinatione grad. 15. m. 43. Sumo ex tabula benefica numerum 103879 quem multiplico in 96134. sinum scilicet secundum arcus graduum 15. m. 59. quod fuit argumentum declinationis: Et produco 9986303786. de quo abicio quinque figuras ad dextram, et relinquunt 99863 cui sinui ascribitur arcus graduum fere 87. qui de quadrante summotus residuat gradus 3 ut prius differentiam scilicet transitus stellæ per coeli medium. Et quoniam locus stellæ fuit in primo Zodiaco quadrante, cum declinatione septemtrionali, idcirco subtraho præfatam differentiam 3. grad. 3. ab arcu tertio graduum 64. m. 47. et supersunt gradus 61. m. 47. ascensio scilicet recta stellæ propositæ quæsita.
\note{« \uncertain{Demo} abicio » : le premier mot se lit mal (« Remoueo » est une lecture du trait) et le second le répète en quelque sorte ; il n'est pas biffé. « dicta » : un trait traverse le mot, sans qu'on puisse dire s'il est biffé. « 9986303786 » : le 6 est ajouté au-dessus de la ligne. « 3. grad. 3. » est écrit ainsi ; le calcul (64 gr. 47 m. − 3 gr. = 61 gr. 47 m.) n'emploie que 3 degrés.}

6. Ad propositam loci latitudinem et declinationem datam puncti cuiuspiam latitudinem ortivam exhibere. Cum latitudine \struck{ortiva} loci ex tabula benefica sume numerum quem multiplica in sinum declinationis data et a producto abice quinque dextrorsum figuras, nam sic relinquetur sinus quæsitæ latitudinis ortivæ. Exempli gratia ad latitudinem Messanæ graduum 38. pro sole habente maximam declinationem graduum 23. m. 30. quæro latitudinem ortivam. Cum gradu 38 accipio ex tabula benefica numerum 126902. quem multiplico in 39875 sinum graduum 23. m. 30. data declinationis et produco 5060217250. Unde abrectis quinque ad dextram figuris relinquitur numerus 50602 qui sinus est quæsitæ latitudinis ortivæ scilicet graduum 30. m. 24.
\note{le « 6 » est en marge, devant l'énoncé : numéro du problème ; les problèmes 1 à 5 sont aux pages de droite de la vue 18 et de gauche de la vue 19. « abrectis » est écrit ainsi. La page est blanche au-dessous.}

\note{page de droite de la vue 18, numérotée à l'encre « 17 » en haut à droite. Elle ne continue pas la page de gauche, qui finit sur le problème 6 : elle commence les problèmes, numérotés en marge 1 à 4.}

1. Puncti Ecclipticæ dati declinationem comperire. Sinum rectum arcus ecclipticæ inter æquinoctium et datum punctum duc in sinum rectum maximæ declinationis, et productum partire in sinum totum: nam sinus ex divisione exeuntis arcus erit declinatio quæsita. Exempli gratia quæro declinationem 11. gradus Taurus cuius finis distat ab æquinoctio per gradus 41 cuius arcus sinus habet partes 65606 quem duco in 39875 sinum graduum 23. m. 30. maximæ declinationis, et produco 2616039250. quod partior in 100.000 sinum totum abiectis scilicet quinque ad dextram figuris, et provenerunt 26160 cuius sinus arcus est \uncertain{15} graduum m. 10 declinatio quæsita.
\note{« exeuntis » et « Taurus » sont écrits ainsi. « 15 » : un 1 et un 5 sous le « g » de « graduum », qui les recouvre en partie ; le problème suivant reprend ce point avec « g. 15. m. 10. ».}

2. Ascensionem rectam dati Ecclipticæ puncti exquirere: Cum declinatione puncti propositi et cum complemento maximæ declinationis sume numeros ex tabula foecunda quos duc \struck{in} invicem et a producto abice quinque ad dextram figuras. Nam sinus relicti arcus erit ascensio recta quæsita. sed ab æquinoctio proximo computanda. Exempli gratia cum g. 15. m. 10. declinatione et puncti prædicti ex tabula foecunda excipio numerum 27108 et cum gradu 66 m. 30. ex eadem habeo numerum 230094 quos multiplico alterutrum et produco 6237388152 hinc abiicio 5 figuras dextras et restant 62374. Cuius sinus arcus est gradus 38. ascensio quæsita. Similiter ex tabula foecunda cum declinatione stellæ propositæ et latitudine regionis eliciatur differentia ascensionalis eiusdem stellæ.
\note{« 230094 » : lecture des six chiffres, dont l'avant-dernier est mal formé ; le produit donné, 6237388152, est celui de 27108 par 230094.}

3. Rectam dati Ecclipticæ puncti ascensionem aliter explorare. Cum declinatione dati puncti sume numerum ex tabula benefica quem duc in sinum secundum arcus sui ecclipticæ ab æquinoctio viciniori, et a producto abice quinque ad dextram figuras, et sinus relicti arcus de quadrante sublatus residuabit ascensionem quæsitam: a viciniori æquinoctij puncto computandam. Exempli gratia cum gradus 15. m. 10 declinatione scilicet puncti prædicti ex tabula benefica elicio numerum 103612 quem duco in 75471 sinum scilicet graduum 41 et produco 7819701252 Unde proiectis quinque ad dextram figuris relinquitur numerus 78197 cuius sinus arcus est graduum 51. m. 26. cuius arcus complementum est graduum 38. m. 34. ascensio quæsita.

4. Stellæ cuiuslibet propositæ declinationem sciscitari. Cum arcu Ecclipticæ inter locum longitudinis stellæ et punctum æquinoctij proximum ingredere tabulam declinationum et ascensionum generalem et accipe tres arcus singulos per gradus et minuta notatos quorum primus est arcus circuli latitudinis \struck{circuli} stellæ tuæ inter æquatorem et Ecclipticam clausus. Secundus est angulus sub dicto latitudinis circulo et æquatore comprehensus. Tertius est arcus æquatoris a puncto proximo æquinoctij usque ad dictum latitudinis
\note{« circuli » est écrit deux fois, la seconde biffée. La page s'arrête sur « latitudinis », au milieu d'une phrase ; le reste de la page est blanc.}

\page{19}
\note{double page ; page de gauche. Elle continue la page de droite de la vue 18 (« usque ad dictum latitudinis » / « circulum receptus »).}
circulum receptus. Horum arcuum primus est septentrionalis, vel meridionalis cuius scilicet nominis est punctum longitudinis stellæ tuæ. Hunc iunge cum latitudine tuæ stellæ, quoniam sunt eiusdem nominis, vel subtrahe minorem de majori, si sunt diversorum nominum. Nam adgregatum eiusdem nominis \struck{cum ambobus} cum ambobus, aut relictum eiusdem nominis cum majori vocabitur argumentum declinationis. Et eiusdem nominis erit declinatio. Estque arcus circuli latitudinis inter æquinoctialem et centrum stellæ qui si nullus esset post subtractionem nullam stella pateretur declinationem, Itaque sinum rectum talis argumenti, multiplica in sinum arcus secundi, qui sinus est ipse numerus multiplicandus quo utitur Iohannes Regiomontanus et productum partire in sinum totum quod facies more solito quinque dexterioribus figuris. Et exibit sinus rectus declinationis quæsitæ et eiusdem nominis cum argumento præfato. Exempli \add{gratia} quæro declinationem stellæ aldebaran quæ oculus Tauri dicitur. cuius longitudo hac tempestate habet gradus 2. m. 50 Geminorum. Latitudo autem meridionalis graduum 5. m. 10. Huius longitudinis punctum ab æquinoctio viciniori hoc est a principio arietis distat g. 62. m. 50. cum quo arcu intro in tabulam declinationum et ascensionum rectarum. Et excipio tres arcus primum grad. 21. m. 5. secundum g. 79. m. 10. \uncertain{vel} arcu primo graduum 21. m. 9. et supersunt gradus 15. m. 59. partis septen\textsuperscript{lis}. quod est argumentum declinationis, cuius sinus est 27536. particularum: in quem multiplico sinum secundi arcus gr. 79. m. 31. scilicet partium 98331. ex tabula cuius sinus totus habet partes 100.000. Et produco 2707642916. De quo abiectis quinque ad dextram figuris relinquentur partes. 27076. cui sinui debetur arcus graduum 15. m. 43. declinatio stellæ quæsita septentrionalis quoniam argumentum septemtrionale fuit.
\note{« cum ambobus » est écrit deux fois, la première fois biffé. « gratia » n'est pas écrit après « Exempli ». Le mot rendu « vel », abrégé « ul », se lit mal. Les nombres sont tels qu'écrits : le second arc est « g. 79. m. 10. », puis « gr. 79. m. 31. » ; le premier « grad. 21. m. 5. », puis « graduum 21. m. 9. », et 21 gr. 9 m. moins 5 gr. 10 m. donne 15 gr. 59 m.}

5. Rectam stellæ cuiusuis ascensionem elucubrare. Primum cum longitudine et latitudine stellæ tuæ supputa modo prædicto eius declinationem per dictam tabulam et prædictos tres arcus. Atque cum declinatione et complemento arcus secundi excipe ex tabula foecunda numeros quos invicem multiplica: et a producto quinque dextrorsum figuras abice, relinquetur enim sinus cuiusdam arcus æquatoris qui circulis latitudinis et declinationis stellæ interiacet. Atque differentia transitus stellæ per coeli medium vocari solet. Vel sic. Cum declinatione dictæ stellæ quære numerum ex tabula benefica quem multiplica in sinum secundum argumenti declinationis, et productum partire in sinum totum, quod facies abiectis quinque dextris figuris et habebis sinum cuius arcus de quadrante ablatus residuabit dictam differentiam transitus stellæ per coeli medium. Itaque ex differentia ista arcuque tertio superius invento in tabula declinationum elicies rectam stellæ ascensionem. Nam in primis
\note{la phrase s'arrête sur « Nam in primis » au bas du texte, et la suite est en tête de la vue 17 à droite, feuillet numéroté 16 : « considerandum est in quo ex quatuor quadrantibus Zodiaci sit stella locus ». Le texte des problèmes se lit donc sur les pages numérotées 17 (vue 18, droite), puis vue 19, gauche, puis 16 (vue 17, droite), puis vue 18, gauche, où il finit avec le problème 6 : les deux feuillets numérotés 16 et 17 sont reliés dans l'ordre inverse de celui du texte. C'est une reconstitution de cette transcription ; les numéros à l'encre, 16 puis 17, se suivent dans l'ordre présent, que le fichier garde. Sur le bord droit de la vue, près du pli, on voit la colonne des degrés (1 à 30 et plus) de la table de la page d'en face.}

\section*{Tabella benefica}
\note{page de droite de la vue 19, numérotée à l'encre « 18 » en haut à droite, occupé par une seule table en deux blocs, degrés 1 à 45 et 46 à 90 ; colonnes intitulées en tête « g. » (« gr. » au second bloc), « Radius », « Dria » (abrégé pour « Differentia »), « 60\textsuperscript{ma} ». Donnée ici en trois tranches de quinze lignes. On ne reproduit ni les points qui précèdent la plupart des différences (« .. 15 », « · 2534 ») ni ceux qui séparent parfois les milliers (« 178·830 »). Aux lignes 4 et 5 de la dernière colonne, « 2.47 » et « 2:18 » sont serrés sur une seule ligne. Au second bloc, la ligne 90 porte « Infinitum » dans les trois colonnes.}

\[
\begin{array}{r|r|r|r||r|r|r|r}
\text{g.} & \text{Radius} & \text{Dria} & 60^{\text{ma}} & \text{gr.} & \text{Radius} & \text{Dria} & 60^{\text{ma}} \\
\hline
1 & 100015 & 15 & 0{:}15 & 46 & 143955 & 2534 & 42.14 \\
2 & 100061 & 46 & 0.46 & 47 & 146628 & 2673 & 44.33 \\
3 & 100137 & 76 & 1.16 & 48 & 149448 & 2820 & 47.0 \\
4 & 100244 & 107 & 2.47 & 49 & 152425 & 2977 & 49.37 \\
5 & 100382 & 138 & 2{:}18 & 50 & 155572 & 3147 & 52.27 \\
\hline
6 & 100551 & 169 & 2{:}49 & 51 & 158902 & 3330 & 42.14 \\
7 & 100751 & 200 & 3.20 & 52 & 162427 & 3525 & 44.33 \\
8 & 100983 & 232 & 3.52 & 53 & 166165 & 3738 & 47.0 \\
9 & 101246 & 264 & 4.29 & 54 & 170131 & 3966 & 49.37 \\
10 & 101543 & 297 & 4{:}57 & 55 & 174344 & 4213 & 52.27 \\
\hline
11 & 101871 & 328 & 5.28 & 56 & 178830 & 4486 & 55.30 \\
12 & 102234 & 363 & 6.3 & 57 & 183608 & 4778 & 58.45 \\
13 & 102630 & 396 & 6.36 & 58 & 188708 & 5100 & 62.18 \\
14 & 103061 & 431 & 7.11 & 59 & 194160 & 5452 & 66.6 \\
15 & 103528 & 467 & 7.47 & 60 & 200000 & 5840 & 70.13 \\
\end{array}
\]

\[
\begin{array}{r|r|r|r||r|r|r|r}
16 & 104630 & 502 & 8.22 & 61 & 206267 & 6267 & 74.46 \\
17 & 104569 & 539 & 8.59 & 62 & 213006 & 6739 & 79.38 \\
18 & 105146 & 577 & 9.37 & 63 & 220269 & 7263 & 85.0 \\
19 & 105762 & 616 & 10.16 & 64 & 228117 & 7848 & 90.52 \\
20 & 106418 & 656 & 10.56 & 65 & 236620 & 8503 & 97.20 \\
\hline
21 & 107115 & 697 & 11.37 & 66 & 245859 & 9239 & 104.27 \\
22 & 107854 & 739 & 12.19 & 67 & 255930 & 10071 & 112.19 \\
23 & 108636 & 782 & 13.2 & 68 & 266947 & 11017 & 121.3 \\
24 & 108464 & 828 & 13.48 & 69 & 279043 & 12096 & 130.48 \\
25 & 110338 & 874 & 14.34 & 70 & 292380 & 13337 & 141.43 \\
\hline
26 & 111260 & 922 & 15.22 & 71 & 307155 & 14775 & 153.59 \\
27 & 112233 & 973 & 16.13 & 72 & 323607 & 16452 & 167.51 \\
28 & 113257 & 1024 & 17.4 & 73 & 342030 & 18423 & 183.57 \\
29 & 114335 & 1078 & 17.58 & 74 & 362796 & 20766 & 201.36 \\
30 & 115470 & 1130 & 18.55 & 75 & 386370 & 23574 & 222.17 \\
\end{array}
\]

\[
\begin{array}{r|r|r|r||r|r|r|r}
31 & 116664 & 1194 & 19.54 & 76 & 413357 & 26987 & 246.15 \\
32 & 117918 & 1254 & 20.54 & 77 & 444541 & 31184 & 274.12 \\
33 & 119236 & 1318 & 21.58 & 78 & 480973 & 36432 & 307.3 \\
34 & 120621 & 1385 & 23.5 & 79 & 524084 & 43111 & 346.6 \\
35 & 122078 & 1457 & 24.17 & 80 & 575877 & 51793 & 392.54 \\
\hline
36 & 123606 & 1528 & 25.28 & 81 & 639245 & 63368 & 449.47 \\
37 & 125214 & 1608 & 26.48 & 82 & 718530 & 79285 & 519.44 \\
38 & 126902 & 1688 & 28.8 & 83 & 820552 & 102022 & 607.12 \\
39 & 128676 & 1774 & 29.34 & 84 & 956677 & 136125 & 718.31 \\
40 & 130541 & 1865 & 31.5 & 85 & 1147371 & 190694 & 863.14 \\
\hline
41 & 132501 & 1960 & 32.40 & 86 & 1433558 & 286187 & 4769.47 \\
42 & 134563 & 2062 & 34.22 & 87 & 1910732 & 477174 & 7952.54 \\
43 & 136733 & 2170 & 36.10 & 88 & 2865371 & 954639 & 15910.39 \\
44 & 139016 & 2283 & 38.3 & 89 & 5729868 & 2864497 & 47741.37 \\
45 & 141421 & 2405 & 40.5 & 90. & \text{Infinitum} & \text{Infinitum} & \text{Infinitum} \\
\end{array}
\]

\note{tel qu'écrit, et que la présente lecture a contrôlé en ajoutant chaque différence au rayon précédent : la colonne des rayons et celle des différences s'accordent partout, sauf lignes 16 (« 104630 », que la somme donnerait 104030), 24 (« 108464 », chiffre surchargé ; la somme donnerait 109464) et 30 (115470 et 1130, la somme donnerait 115465) ; ligne 25, « 110338 », le 3 est surchargé ; ligne 44, un trait de trop après « 139016 » ; ligne 43 de la dernière colonne, « 36.10 », le 6 formé comme un b. Dans la dernière colonne du second bloc, les lignes 51 à 55 répètent les valeurs des lignes 46 à 50, et de la ligne 56 à la ligne 85 chaque valeur est le soixantième de la différence écrite cinq lignes plus haut ; les lignes 86 à 89 s'accordent de nouveau avec leur propre différence. Ligne 4 de la même colonne, « 2.47 » là où la différence 107 donne 1:47. Ligne 86, « 4769.47 » : le quatrième chiffre se lit aussi 4. Ligne 43, « 36.10 » se lit aussi « 39.20 ». Ligne 89, « 47741.37 » : le cinquième chiffre, surchargé, se lit aussi 5.}

\[
\begin{array}{lcl}
0.\;\;0. & \ldots\ldots & 100.000 \\
89.\;15 & \ldots\ldots & 7639653 \\
89.\;30 & \ldots\ldots & 11459309 \\
89.\;45 & \ldots\ldots & 22918381 \\
89.\;55 & \ldots\ldots & 68754512 \\
89.\;59 & \ldots\ldots & 343772560.
\end{array}
\]
\note{ces six lignes sont écrites sous la table, comme sous la « Tabula foecunda » de la vue 17 ; la page est blanche au-dessous}

\page{20}
\note{double page ; page de gauche, une table sans titre, en deux blocs, degrés 1 à 45 et 46 à 90 ; pour chaque degré trois arcs, « Primus », « Secundus », « Tertius », en degrés et minutes (« g. », « gr. », « m. »). C'est, par ses trois arcs, la « tabula declinationum et ascensionum » que les problèmes 4 et 5 (vues 18 et 19) disent employer ; le rapprochement est de cette lecture. Donnée ici en six tranches de quinze lignes, les deux blocs l'un après l'autre. La moitié inférieure de la page est blanche.}

\[
\begin{array}{r|rr|rr|rr}
& \text{Primus} & & \text{Secundus} & & \text{Tertius} & \\
\text{g.} & \text{gr.} & \text{m.} & \text{g.} & \text{m.} & \text{g.} & \text{m.} \\
\hline
1 & 0 & 26 & 66 & 30 & 1 & 6 \\
2 & 0 & 52 & 66 & 31 & 2 & 11 \\
3 & 1 & 18 & 66 & 32 & 3 & 16 \\
4 & 1 & 44 & 66 & 33 & 4 & 22 \\
5 & 2 & 20 & 66 & 36 & 5 & 27 \\
\hline
6 & 2 & 36 & 66 & 38 & 6 & 32 \\
7 & 3 & 2 & 66 & 41 & 7 & 38 \\
8 & 3 & 28 & 66 & 44 & 8 & 43 \\
9 & 3 & 53 & 66 & 48 & 9 & 48 \\
10 & 4 & 19 & 66 & 52 & 10 & 52 \\
\hline
11 & 4 & 45 & 66 & 57 & 11 & 58 \\
12 & 5 & 10 & 67 & 2 & 13 & 3 \\
13 & 5 & 35 & 67 & 8 & 14 & 8 \\
14 & 6 & 0 & 67 & 14 & 15 & 13 \\
15 & 6 & 25 & 67 & 21 & 16 & 17 \\
\end{array}
\]

\[
\begin{array}{r|rr|rr|rr}
16 & 6 & 50 & 67 & 28 & 17 & 22 \\
17 & 7 & 15 & 67 & \uncertain{35} & 18 & 27 \\
18 & 7 & 39 & 67 & \uncertain{43} & 19 & 31 \\
19 & 8 & 3 & 67 & \uncertain{51} & 20 & 35 \\
20 & 8 & 27 & 67 & \uncertain{59} & 21 & 39 \\
\hline
21 & 8 & 51 & 68 & 8 & 22 & 43 \\
22 & 9 & 15 & 68 & 18 & 23 & 47 \\
23 & 9 & 39 & 68 & 28 & 24 & 51 \\
24 & 10 & 2 & 68 & 38 & 25 & 54 \\
25 & 10 & 25 & 68 & 49 & 26 & 57 \\
\hline
26 & 10 & 48 & 69 & 0 & 28 & 0 \\
27 & 11 & 10 & 69 & 11 & 29 & 3 \\
28 & 11 & 32 & 69 & 23 & 30 & 6 \\
29 & 11 & 54 & 69 & 35 & 31 & 9 \\
30 & 12 & 16 & 69 & 48 & 32 & 11 \\
\end{array}
\]

\[
\begin{array}{r|rr|rr|rr}
31 & 12 & 37 & 70 & 1 & 33 & 14 \\
32 & 12 & 58 & 70 & 14 & 34 & 16 \\
33 & 13 & 19 & 70 & 28 & 35 & 18 \\
34 & 13 & 40 & 70 & 42 & 36 & 20 \\
35 & 14 & 0 & 70 & 56 & 37 & 22 \\
\hline
36 & 14 & 20 & 71 & 11 & 38 & 23 \\
37 & 14 & 30 & 71 & 26 & 39 & 25 \\
38 & 14 & 59 & 71 & 41 & 40 & 26 \\
39 & 15 & 18 & 71 & 57 & 41 & 27 \\
40 & 15 & 37 & 72 & 13 & 42 & 28 \\
\hline
41 & 15 & 55 & 72 & 29 & 43 & 28 \\
42 & 16 & 13 & 72 & 45 & 44 & 28 \\
43 & 16 & 31 & 73 & 2 & 45 & 29 \\
44 & 16 & 48 & 73 & 19 & 46 & 29 \\
45 & 17 & 5 & 73 & 37 & 47 & 29 \\
\end{array}
\]

\[
\begin{array}{r|rr|rr|rr}
& \text{Primus} & & \text{Secundus} & & \text{Tertius} & \\
\text{gr.} & \text{g.} & \text{m.} & \text{g.} & \text{m.} & \text{g.} & \text{m.} \\
\hline
46 & 17 & 22 & 73 & 55 & 48 & 29 \\
47 & 17 & 38 & 74 & 13 & 49 & 28 \\
48 & 17 & 54 & 74 & 31 & 50 & 27 \\
49 & 18 & 10 & 74 & 50 & 51 & 26 \\
50 & 18 & 25 & 75 & 9 & 52 & 25 \\
\hline
51 & 18 & 40 & 75 & 29 & 53 & 24 \\
52 & 18 & 55 & 75 & 47 & 54 & 23 \\
53 & 19 & 9 & 76 & 8 & 55 & 21 \\
54 & 19 & 23 & 76 & 27 & 56 & 19 \\
55 & 19 & 36 & 76 & 47 & 57 & 18 \\
\hline
56 & 19 & 49 & 77 & 7 & 58 & 16 \\
57 & 20 & 2 & 77 & 27 & 59 & 14 \\
58 & 20 & 14 & 77 & 48 & 60 & 12 \\
59 & 20 & 26 & 78 & 9 & 61 & 9 \\
60 & 20 & 38 & 78 & 30 & 62 & 6 \\
\end{array}
\]

\[
\begin{array}{r|rr|rr|rr}
61 & 20 & 49 & 78 & 51 & 63 & 3 \\
62 & 21 & 0 & 79 & 13 & 64 & 0 \\
63 & 21 & 11 & 79 & 34 & 64 & 57 \\
64 & 21 & 21 & 79 & 56 & 65 & 54 \\
65 & 21 & 31 & 80 & 18 & 66 & 51 \\
\hline
66 & 21 & 40 & 80 & 40 & 67 & 47 \\
67 & 21 & 49 & 81 & 2 & 68 & 44 \\
68 & 21 & 58 & 81 & 22 & 69 & 40 \\
69 & 22 & 6 & 81 & 47 & 70 & 36 \\
70 & 22 & 14 & 82 & 20 & 71 & 33 \\
\hline
71 & 22 & 21 & 82 & 32 & 72 & 29 \\
72 & 22 & 28 & 82 & 55 & 73 & 25 \\
73 & 22 & 35 & 83 & 18 & 74 & 21 \\
74 & 22 & 41 & 83 & 41 & 75 & 17 \\
75 & 22 & 47 & 84 & 4 & 76 & 12 \\
\end{array}
\]

\[
\begin{array}{r|rr|rr|rr}
76 & 22 & 52 & 84 & 27 & 77 & 8 \\
77 & 22 & 57 & 84 & 50 & 78 & 3 \\
78 & 23 & 2 & 85 & 14 & 78 & 58 \\
79 & 23 & 7 & 85 & 37 & 79 & 54 \\
80 & 23 & 11 & 86 & 1 & 80 & 49 \\
\hline
81 & 23 & 15 & 86 & 24 & 81 & 44 \\
82 & 23 & 18 & 86 & 48 & 82 & 40 \\
83 & 23 & 21 & 87 & 12 & 83 & 35 \\
84 & 23 & 23 & 87 & 37 & 84 & 30 \\
85 & 23 & 25 & 88 & 1 & 85 & 25 \\
\hline
86 & 23 & 27 & 88 & 25 & 86 & 20 \\
87 & 23 & 28 & 88 & 48 & 87 & 15 \\
88 & 23 & 29 & 89 & 12 & 88 & 10 \\
89 & 23 & 30 & 89 & 33 & 89 & 5 \\
90 & 23 & 30 & 90 & 0 & 90 & 0 \\
\end{array}
\]

\note{lignes 17 à 20, colonne « Secundus », minutes : chiffres surchargés où l'on voit deux tracés à la fois (35 et 38 ; 43 et 48 ; 51 et 54 ; 50 et 59) ; la lecture offerte est celle qui suit la progression de la colonne (28, puis 68 gr. 8 m. à la ligne 21), ce qui est une inférence. Ligne 70, « Secundus », « 82. 20 » tel qu'écrit, entre 81. 47 et 82. 32. Ligne 79, « Tertius », « 79. 54 » : le dernier chiffre se lit aussi 9. Tel qu'écrit : ligne 5, « Primus », « 2. 20 » (le 2 des minutes est précédé d'un point et se lit peut-être 1) ; ligne 18, « 7. 39 », le dernier chiffre se lit aussi 4 ; ligne 37, « 14. 30 », entre 14. 20 et 14. 59. Chiffres surchargés mais lisibles : lignes 16 (« Tertius », 22), 27 (« Primus », 11), 43 et 83 (« Secundus », 73 et 87), 53 (« Primus », 19), 63 (« Secundus », 34).}

\note{page de droite de la vue 20, numérotée à l'encre « 19 » en haut à droite. Elle ne continue pas la page de gauche, une table : elle s'ouvre sur la fin de la démonstration d'une proposition dont l'énoncé et le début ne sont pas dans ce lot. Les propositions XXIII et XXIIII portent sur les triangles rectilignes, la XXV sur le triangle sphérique rectangle ; la vue 21 (lot 2) continue la XXV (« Sit triangulum sphærale abc »). Dans la marge droite, en face de la proposition XXIIII, le cachet rond « Bibliothèque impériale MSS ».}

et per eamdem et conversam proportionem sicut sinus totus ad sinum anguli $cbd$ sic sinus secundus arcus $bd$ ad sinum secundum anguli $c$. Igitur ex æquali sicut sinus anguli $abd$ ad sinum anguli $cbd$ sic sinus secundus anguli $a$ ad sinum secundum anguli $c$. Quod est propositum.

\subsection*{Propositio XXIII.}

Duo triangula rectilinea invicem æqualia et æquiangula inter se sunt æquilatera.

Sunto duo triangula rectilinea $abc$ $def$ invicem æqualia et æquiangula hoc est angulos $ad$ æquales invicem et angulos $be$ æquales: et angulos $cf$ æquales invicem habentia. Aio quod et inter se sunt æquilatera hoc est latera æquis angulis opposita æqualia invicem habebunt. Secus enim ponatur ipsi triangulo $def$ æquilaterum triangulum $agh$. eruntque per 4. 6. et 4. 1. triangula $def$, $agh$ æqualia, et idcirco triangula $abc$ $agh$ æqualia pars et totum quod est impossibile. Astruitur ergo propositum.
\note{en marge à droite, deux figures : un triangle lettré $a$, $b$, $c$ et un triangle lettré $d$, $e$, $f$. Dans le texte, les lettres qui nomment points, angles et arcs sont souvent surmontées d'un trait, qu'on ne reproduit pas.}

\subsection*{Propositio XXIIII.}

Duo triangula rectilinea invicem æqualia eamdem sive æquas bases habentia et basibus oppositos angulos æquales habebunt et reliqua latera singula singulis æqualia.

Sunto duo triangula Rectilinea $abc$ $eda$ super eamdem basem ac invicem æqualia et angulos $bd$ æquales habentia. Aio quod eorum duo latera $ab$ $ed$ erunt æqualia. Itemque duo latera reliqua $bc$ $da$ æqualia erunt. Talia enim triangula circumscribuntur ab eodem circulo et sunt inter lineas æquidistantes et æqualia et æquiangula et ideo per præcedentem æquilatera. quod demonstrandum proponebatur.
\note{en marge à droite, une figure : un quadrilatère lettré $d$, $b$ en haut, $a$, $c$ en bas, dont les diagonales se coupent en $e$. « æqualia » est écrit deux fois, en fin de ligne puis au début de la suivante (« ab ed erunt æqualia æqualia »), et n'est donné ici qu'une fois.}

\subsection*{Propositio XXV.}

Si trianguli rectanguli ex arcubus circulorum majorum in superficie Sphæræ et quadrantibus minorum sinus unus secundus arcuum acutis oppositorum fuerit medius proportionalis inter sinum totum sinumque acuti anguli oppositi, tunc reliqui arcus erunt alter complemento alterius æquales et sinus secundus prædicti arcus acutum subtendentis æqualis erit sinui reliqui acuti anguli.
\note{« secundus » est écrit dans l'interligne au-dessus de « unus », qui n'est pas biffé ; les deux mots sont donnés dans l'ordre où ils se lisent. La page s'arrête sur « anguli », le reste est blanc ; l'énoncé est démontré à la vue 21 (lot 2).}

\end{document}
